\documentclass[11pt]{article}

\usepackage[margin=1.05in]{geometry}
\usepackage{amsmath,amssymb,amsthm}
\usepackage{graphicx}
\usepackage{booktabs}
\usepackage{xcolor}
\usepackage[colorlinks=true,linkcolor=blue!55!black,citecolor=blue!55!black,
            urlcolor=blue!55!black]{hyperref}

\newtheorem{proposition}{Proposition}
\newtheorem{corollary}[proposition]{Corollary}
\theoremstyle{definition}
\newtheorem{remark}[proposition]{Remark}
\newtheorem{recipe}[proposition]{Recipe}

\newcommand{\abar}{\bar\alpha}
\newcommand{\snr}{\mathrm{SNR}}
\newcommand{\sig}{\tilde\sigma}
\newcommand{\Var}{\mathrm{Var}}
\newcommand{\dd}{\mathrm{d}}
\newcommand{\R}{\mathbb{R}}
\newcommand{\E}{\mathbb{E}}
\newcommand{\Tr}{\mathrm{tr}}
\newcommand{\mmse}{\mathrm{mmse}}
\newcommand{\leff}{\lambda^{\mathrm{eff}}}

\title{\bf The noise schedule should be specified in effective log-SNR\\
\large What the cosine schedule assumes about the data, and how to remove it}
\author{}
\date{\today}

\begin{document}
\maketitle
\vspace{-2.4em}

\begin{abstract}
\noindent
A diffusion schedule is a plan for spending $T$ steps along the noise axis. We
show that every criterion one might use to lay out those steps --- information
delivered, or variance transported by the sampler --- is a function of the
\emph{effective} log-SNR $\leff=\lambda+2\log s$, where $s$ is the standard
deviation of the clean data, and never of the nominal $\lambda$ alone. The
information statement is exact and distribution-free: by the I-MMSE relation,
$\dd I/\dd\lambda=\tfrac12 e^{\leff}\mmse(e^{\leff})$, with no high-SNR
approximation anywhere. Since a schedule controls $\lambda$ while the process
responds to $\leff$, and the two differ by the constant $2\log s$, any
schedule specified purely in terms of $\abar$ has silently assumed $s=1$. The
correction is a rigid translation of the log-SNR curve by $2\log(1/\hat s)$,
which is two lines of code and is \emph{exactly} equivalent to the
reparametrization by polar angle in the signal--noise plane --- the latter is a
derived consequence, not a premise. We compute the exact per-step DDIM variance
deficit $s^2\sin^2(\Delta\theta_k)/v_k$, obtain $D_{\min}=\pi^2/(4T)$ and
$D=\pi^2/(4T)\cdot\cosh\log(s/\hat s)$, and show the unshifted cosine schedule is
the special case $\hat s=1$. For heterogeneous or state-dependent $s$, two
scalars suffice: $M_+=\E[s]$ and $M_-=\E[1/s]$ give the best global shift
$\hat s=\sqrt{M_+/M_-}$ and the irreducible penalty $\sqrt{M_+M_-}$. Because the
penalty sees only the log-error it cannot distinguish spread from estimation
noise, so the same rule applied to a posterior turns the cost into a loss: from
$K$ samples per state the optimal single-state shift \emph{inflates} the sample
width to $\hat s_K\sqrt{\nu/(\nu-1)}$, $\nu=K-1$, at residual cost
$\sqrt{f(\nu)}\approx1+1/(4\nu)$. Having identified the log-error as what costs,
the estimator itself needs no further decision theory: a global shift is obtained
by averaging $\log\hat s_K$ and removing the exact additive bias
$\tfrac12(\psi(\nu/2)-\log(\nu/2))$, whose second moment simultaneously supplies
the heterogeneity floor. That targets the geometric mean, and since the risk in
$\log\hat s$ is itself exactly a $\cosh$, the substitution costs precisely
$\cosh$ of its log-distance to $\sqrt{M_+/M_-}$ --- zero when $\log s$ is
symmetric. Nothing depends on the mean. All closed forms are verified
numerically.
\end{abstract}

\section{The claim, and the fix}

A variance-preserving diffusion is usually specified by a sequence
$\abar_k\in(0,1)$, equivalently by a trajectory of log-SNR values
$\lambda_k=\log[\abar_k/(1-\abar_k)]$. The cosine schedule
\cite{nichol2021improved} is one such trajectory and it is close to universal in
practice.

This note makes one point. The quantity that a schedule writes down is
$\lambda_k$, but the quantity that the process responds to is the log-SNR
\emph{of the actual data},
\begin{equation}
  \leff_k \;=\; \lambda_k + 2\log s,
  \label{eq:shift0}
\end{equation}
with $s$ the standard deviation of the clean distribution. Sections
\ref{sec:info} and \ref{sec:transport} show that both natural design criteria
--- how much information the channel delivers, and how much variance the sampler
fails to transport --- are functions of $\leff$ alone. A trajectory written in
$\lambda$ is therefore a claim about $\leff$ only when $s=1$. Since \eqref{eq:shift0}
is a constant offset, the repair is a rigid vertical translation:

\begin{recipe}[shifted schedule]
\label{rec:main}
Let $\lambda^{\rm target}_k$ be any trajectory you trust for unit-variance data
(e.g.\ cosine's), and let $\hat s$ estimate the clean standard deviation. Use
\begin{equation}
  \boxed{\;
  \lambda_k \;=\; \lambda^{\rm target}_k - 2\log\hat s,
  \qquad
  \abar_k \;=\; \frac{1}{1+e^{-\lambda_k}}
  \;}
  \label{eq:recipe}
\end{equation}
equivalently $\sig_k=\hat s\,\sig^{\rm target}_k$, equivalently
$\abar_k=1/(1+\hat s^2\,\sig^{{\rm target}\,2}_k)$, where
$\sig^2=(1-\abar)/\abar$. Then $\leff_k=\lambda^{\rm target}_k$ exactly.
\end{recipe}

Everything else here is either a justification of \eqref{eq:recipe}, a
computation of what it buys, or a statement of what to do when a single $\hat s$
is not enough. Three facts about \eqref{eq:recipe} are worth stating up front
because they are easy to get wrong:

\begin{itemize}\itemsep2pt
  \item \textbf{It is exact, not approximate.} Applying \eqref{eq:recipe} to the
  cosine trajectory produces precisely the schedule whose steps are uniform in
  the polar angle $\psi=\arctan(\sig/s)$ of the signal--noise plane
  (Proposition~\ref{prop:equiv}); verified to $10^{-14}$ in
  Section~\ref{sec:verify}. The angle description and the shift description are
  the same object.
  \item \textbf{It is not a reshaping.} A translation in $\lambda$ is not in the
  family swept out by cosine's shape parameters, which is why one cannot obtain
  it by retuning them \cite{chen2023importance}. This is a statement about that
  family, not an impossibility: \eqref{eq:recipe} implements the translation
  directly, in two lines.
  \item \textbf{It is equivalent to standardizing the data.} Diffusing $x_0/\hat s$
  under $\lambda^{\rm target}$ and diffusing $x_0$ under \eqref{eq:recipe} give
  the same geometry. Which one to use is an engineering choice
  (Section~\ref{sec:beyond}); for anisotropic data only standardization can act
  per coordinate.
\end{itemize}

\section{Setup: nominal versus effective log-SNR}
\label{sec:setup}

Fix a clean distribution $p_0$ on $\R^d$ with mean $\mu$ and covariance
$\Sigma$. The forward process is
\begin{equation}
  x_k \;=\; \sqrt{\abar_k}\,x_0 \;+\; \sqrt{1-\abar_k}\,\varepsilon,
  \qquad \varepsilon\sim\mathcal N(0,I),
  \label{eq:forward}
\end{equation}
with $\abar_k$ decreasing in $k$; we index so that \textbf{$k=0$ is cleanest},
matching the code, so the reverse process steps $k\to k-1$. Since $\varepsilon$
is centered and isotropic and \eqref{eq:forward} is affine,
\begin{equation}
  \E[x_k]=\sqrt{\abar_k}\,\mu,
  \qquad
  V_k \;=\; \abar_k\Sigma + (1-\abar_k)I .
  \label{eq:marginal}
\end{equation}
Note that $V_k$ contains no $\mu$; Section~\ref{sec:mean} draws the consequence.

We use four equivalent coordinates for a noise level, and one derived one:
\begin{equation}
  \abar\in(0,1),
  \quad
  \sig=\sqrt{\tfrac{1-\abar}{\abar}},
  \quad
  \lambda=\log\tfrac{\abar}{1-\abar}=-2\log\sig,
  \quad
  \theta=\arctan\sig,
  \quad
  \psi=\arctan\tfrac{\sig}{s}.
  \label{eq:coords}
\end{equation}
Thus $\abar=\cos^2\theta$, with $\theta=0$ clean and $\theta=\pi/2$ pure noise.

\paragraph{The two SNRs.} Take $d=1$ and $\Sigma=s^2$ for the moment. The
schedule's nominal SNR is $\snr=\abar/(1-\abar)=e^{\lambda}$, which is what
\eqref{eq:forward} would deliver if the clean variance were $1$. The ratio the
process actually realizes is
\begin{equation}
  \snr^{\rm eff}_k
  \;=\; \frac{\Var(\text{signal part of }x_k)}{\Var(\text{noise part})}
  \;=\; \frac{\abar_k s^2}{1-\abar_k}
  \;=\; s^2\,e^{\lambda_k},
  \qquad\text{so}\qquad
  \leff_k=\lambda_k+2\log s .
  \label{eq:shift}
\end{equation}
Writing $v_k=\abar_ks^2+(1-\abar_k)$ for the marginal variance, note
$\cos^2\psi_k=s^2\abar_k/v_k$ is the signal's share of it, so $\psi$ is just
$\leff$ on a compact scale: $\leff=-2\log\tan\psi$. These are three names for
one thing.

The rest of this section is the only assumption we need to make explicit.
\emph{A schedule chooses $\lambda_k$. If the criterion by which we judge a
schedule is a function of $\leff$, then by \eqref{eq:shift} the schedule is
under-determined by one constant until $s$ is named, and setting that constant
by ignoring $s$ amounts to asserting $s=1$.} Sections \ref{sec:info} and
\ref{sec:transport} establish the antecedent, twice, from independent premises.

\begin{remark}[a notation collision]
The Nichol--Dhariwal cosine schedule contains a small offset conventionally
also written $s$ (value $0.008$), whose only job is to keep $\beta_1$ away from
$0$. Here $s$ and $s_i$ always denote clean-data standard deviations; we write
$\delta$ for that offset.
\end{remark}

\section{The information clock, exactly}
\label{sec:info}

The usual argument for log-SNR is that equal steps in $\lambda$ deliver equal
information, justified through the high-SNR limit or through the diffusion
ELBO. Neither is necessary, and the high-SNR version understates the case: the
information rate has an exact closed form at every noise level, for every clean
distribution.

\subsection{The exact rate}

Dividing \eqref{eq:forward} by $\sqrt{1-\abar_k}$ is invertible, so
\begin{equation}
  I(x_0;x_k) \;=\; I\big(x_0;\ \sqrt{\gamma_k}\,x_0+\varepsilon\big),
  \qquad \gamma_k=\frac{\abar_k}{1-\abar_k}=e^{\lambda_k}.
  \label{eq:channel}
\end{equation}
Level $k$ of a diffusion \emph{is} a Gaussian channel of SNR $\gamma_k$. For
that channel the I-MMSE relation of Guo, Shamai and Verd\'u
\cite{guo2005mutual} states that, in nats,
\begin{equation}
  \frac{\dd}{\dd\gamma}I(x_0;\sqrt\gamma x_0+\varepsilon)
  \;=\; \tfrac12\,\mmse(\gamma),
  \qquad
  \mmse(\gamma)=\E\big\|x_0-\E[x_0\mid \sqrt\gamma x_0+\varepsilon]\big\|^2,
  \label{eq:immse}
\end{equation}
for \emph{any} $x_0$ with finite variance. This is an identity, not an
expansion. Since $\dd\gamma=\gamma\,\dd\lambda$,
\begin{equation}
  \boxed{\;
  \frac{\dd I}{\dd\lambda}
  \;=\; \underbrace{\frac{\gamma}{2}}_{\substack{\text{channel}\\\text{quality}}}
        \times
        \underbrace{\mmse(\gamma)}_{\substack{\text{how much is}\\\text{still unknown}}}
  \;}
  \label{eq:rate}
\end{equation}
Equation \eqref{eq:rate} is the precise form of the intuition that two effects
compete. At low SNR there is a great deal left to learn --- $\mmse\to\Var(x_0)$,
the whole prior variance --- but the channel is so poor that almost none of it
arrives per unit $\lambda$. At high SNR the channel is excellent but there is
almost nothing left: $\mmse\sim1/\gamma$, and the two effects cancel exactly,
leaving a constant rate. Neither factor alone is the right clock; their product
is.

\subsection{The clock depends on the data only through $\leff$ and a shape}

\begin{proposition}
\label{prop:clock}
Write $x_0=\mu+sZ$ with $Z$ standardized ($\E Z=0$, $\Var Z=1$) of shape $p_Z$.
Then
\begin{equation}
  \frac{\dd I}{\dd\lambda}
  \;=\; \tfrac12\,R_{p_Z}\!\big(\leff\big),
  \qquad
  R_{p_Z}(u)\;=\;e^{u}\,\mmse_Z(e^{u}),
  \qquad \leff=\lambda+2\log s .
  \label{eq:clockdep}
\end{equation}
The rate depends on the schedule and the data scale \textbf{only through
$\leff$}, and on the data otherwise only through its shape. Moreover
\begin{equation}
  R_{p_Z}(u)\;\xrightarrow[u\to-\infty]{}\;e^{u},
  \qquad
  R_{p_Z}(u)\;\xrightarrow[u\to+\infty]{}\;1
  \quad\text{for every continuous }p_Z,
  \label{eq:asympt}
\end{equation}
and for the Gaussian shape $R(u)=\sigma(u)=(1+e^{-u})^{-1}$ exactly, which
dominates every other shape pointwise.
\end{proposition}

\begin{proof}
With $x_0=\mu+sZ$, the channel $\sqrt\gamma x_0+\varepsilon$ has the same
posterior structure as $\sqrt{\gamma s^2}Z+\varepsilon$ after the invertible
shift and scaling, so $\mmse_{x_0}(\gamma)=s^2\mmse_Z(\gamma s^2)$; the mean
$\mu$ drops out because it is known. Then
$\tfrac\gamma2\mmse_{x_0}(\gamma)=\tfrac12(\gamma s^2)\mmse_Z(\gamma s^2)$, and
$\gamma s^2=e^{\leff}$. As $u\to-\infty$, $\mmse_Z\to\Var Z=1$, giving $e^u$; as
$u\to+\infty$ the posterior of a continuous $Z$ is asymptotically Gaussian with
variance $e^{-u}$, giving $R\to1$. For Gaussian $Z$, $\mmse_Z(\gamma')=1/(1+\gamma')$
so $R(u)=e^u/(1+e^u)$. Pointwise dominance is the standard fact that among all
$Z$ of unit variance the Gaussian has the largest MMSE at every SNR.
\end{proof}

\begin{figure}[t]
  \centering
  \includegraphics[width=\textwidth]{figs/fig_information.pdf}
  \caption{\textbf{The exact information clock.} (a) $\dd I/\dd\leff$ from
  \eqref{eq:clockdep}. For Gaussian data it is exactly $\tfrac12\sigma(\leff)$
  --- no approximation --- rising from the universal low-SNR asymptote
  $\tfrac12e^{\leff}$ to a plateau of $\tfrac12$ nat per unit $\leff$. A
  two-point prior has the same low-SNR asymptote but collapses at high SNR
  because it holds only $\log2$ nats in total. (b) The two competing factors of
  \eqref{eq:rate}: one falls, one rises, and their product saturates. (c)
  Cumulative information. For continuous data it grows without bound at slope
  $\tfrac12$, so no finite schedule can deliver all of it.}
  \label{fig:information}
\end{figure}

Proposition~\ref{prop:clock} is the primal statement this note rests on, and it
already forces Recipe~\ref{rec:main}:

\begin{corollary}[the shift is forced]
\label{cor:forced}
Any schedule criterion built from the information the channel delivers is a
functional of the sequence $\leff_k$ alone. Two schedules with the same
$\{\leff_k\}$ are indistinguishable to it, and a target trajectory expressed in
$\lambda$ realizes the intended $\{\leff_k\}$ if and only if it is translated by
$-2\log s$, which is \eqref{eq:recipe}.
\end{corollary}

Note what this does \emph{not} require: no Gaussian assumption, no high-SNR
limit, and no appeal to the ELBO. It also explains why the correction cannot be
absorbed into a reshaping. Shape changes reparametrize the path
$k\mapsto\lambda_k$; the offset $2\log s$ changes which physical noise levels
those $\lambda_k$ name.

\subsection{Why equal information per step cannot be the whole criterion}

Corollary~\ref{cor:forced} fixes the \emph{coordinate} in which a schedule
should be written. It does not fix the \emph{spacing}, and it cannot, for a
reason visible in Figure~\ref{fig:information}c. Integrating
\eqref{eq:clockdep} for continuous data,
\begin{equation}
  I(\leff) \;=\; \tfrac12\log\big(1+e^{\leff}\big)
  \;\xrightarrow[\leff\to\infty]{}\;\infty
  \qquad\text{at slope }\tfrac12 .
  \label{eq:divergence}
\end{equation}
The total information required to resolve continuous data exactly is infinite,
and it is spread at a constant rate over an infinite stretch of $\leff$. So:

\begin{itemize}\itemsep2pt
  \item \textbf{``Equal information per step'' has no finite-$T$ solution.} Any
  schedule delivering $\Delta I$ per step covers only a bounded range of $\leff$,
  and the endpoints $\abar\to1$, $\abar\to0$ stay infinitely far away
  (Figure~\ref{fig:coordinates}a). To use this criterion one must first truncate
  $\leff$ to a finite window, and the two cutoffs become hyperparameters. That is
  exactly what a uniform-$\lambda$ schedule with an \texttt{snr\_max} cutoff
  does.
  \item \textbf{The divergence is one-sided and structural.} It sits at the
  \emph{clean} end, and only for continuous data: the two-point prior in
  Figure~\ref{fig:information} saturates at $\log2$ nats and needs no cutoff.
  Real data lies between, which is why the cutoff is a modeling choice rather
  than a fact.
\end{itemize}

We therefore need a criterion with a \emph{finite} total, so that minimizing it
at fixed $T$ has a solution. Section~\ref{sec:transport} uses the error the
sampler actually makes. It will return the same coordinate $\leff$ --- which is
the real content of Corollary~\ref{cor:forced} --- and, in addition, a spacing.

\section{What the cosine schedule is, and the identification it makes}
\label{sec:cosine}

The cosine schedule sets $\abar_k=\cos^2\theta_k$ with $\theta_k$ uniform,
\begin{equation}
  \theta_k=\frac{k+1}{T}\cdot\frac{\pi}{2},
  \qquad
  \abar_k=\cos^2\!\Big(\frac{k+1}{T}\frac{\pi}{2}\Big),
  \label{eq:cosine}
\end{equation}
up to the offset $\delta$. Two observations place it relative to
Section~\ref{sec:info}.

\paragraph{It compactifies rather than truncates.} $\theta$ ranges over the
bounded interval $[0,\pi/2]$ (Figure~\ref{fig:coordinates}b), so
\eqref{eq:cosine} reaches the endpoints of the path in finitely many steps with
no cutoff hyperparameter. It is not a truncated uniform-$\lambda$ schedule; it is
a different response to the divergence \eqref{eq:divergence}. Differentiating
$\lambda=-2\log\tan\theta$,
\begin{equation}
  \Big|\frac{\dd\lambda}{\dd\theta}\Big| = \frac{4}{\sin2\theta} \;\geq\; 4,
  \label{eq:speed}
\end{equation}
with equality only at $\theta=\pi/4$, i.e.\ $\snr=1$. So cosine coincides with
uniform-$\lambda$ spacing near $\snr=1$ and takes ever larger log-SNR strides
toward both ends (Figure~\ref{fig:trajectories}).

\paragraph{It is written in $\lambda$, not $\leff$.} Equation \eqref{eq:cosine}
specifies $\abar_k$ and nothing else. By Corollary~\ref{cor:forced} it is a
claim about the process only once $s$ is named, and reading it directly as a
trajectory of $\leff$ sets $s=1$. Everything in the remainder of this note is
the quantification of that step.

\begin{figure}[t]
  \centering
  \includegraphics[width=\textwidth]{figs/fig_coordinates.pdf}
  \caption{\textbf{The unboundedness of $\lambda$ and the compactness of
  $\theta$.} (a) $\lambda$ diverges at both ends of the path. (b)
  $\theta=\arctan\sig$ covers the path in $[0,\pi/2]$. (c) When $s=1$ the point
  (signal amplitude, noise amplitude) traverses the unit circle, on which
  $\theta$ is the polar angle and the arc length. Section~\ref{sec:transport}
  shows what replaces $\theta$ when $s\neq1$.}
  \label{fig:coordinates}
\end{figure}

\begin{figure}[t]
  \centering
  \includegraphics[width=\textwidth]{figs/fig_trajectories.pdf}
  \caption{\textbf{Cosine is not a truncated uniform-$\lambda$ schedule.} (a)
  Log-SNR against step index. Uniform $\lambda$ is a straight line but needs the
  cutoffs $|\lambda|\le12$; cosine is steep at both ends; linear $\beta$ stays
  noisy far too long and then jumps, the pathology cosine was introduced to fix.
  (b) Cosine's log-SNR speed \eqref{eq:speed}, minimized at $\snr=1$.}
  \label{fig:trajectories}
\end{figure}

\section{A criterion that terminates: what the sampler actually loses}
\label{sec:transport}

Section~\ref{sec:info} fixed the coordinate but could not fix the spacing,
because the total information is infinite. We now use a cost that is finite: the
error the reverse sampler makes. Take $p_0=\mathcal N(\mu,\Sigma)$ so the score
is exact and no network error is involved --- by
Proposition~\ref{prop:clock} the Gaussian is also the hardest shape at every
noise level, so this is the conservative choice. Then a deterministic DDIM
sampler makes exactly one kind of error, and we can compute it in closed form.

\subsection{The exact per-step deficit}

\begin{proposition}[exact per-step DDIM deficit]
\label{prop:deficit}
Let $d=1$, $p_0=\mathcal N(m,s^2)$, and let DDIM with $\eta=0$ be fed the exact
$\E[\varepsilon\mid x_k]$. Then the update is affine, $x_{k-1}=c_kx_k+b_km$,
with
\begin{equation}
  c_k=\frac{\sqrt{\abar_{k-1}\abar_k}\,s^2+\sqrt{(1-\abar_{k-1})(1-\abar_k)}}{v_k},
  \qquad v_k=\abar_ks^2+(1-\abar_k),
\end{equation}
it transports the mean \emph{exactly}, $\E[x_{k-1}]=\sqrt{\abar_{k-1}}\,m$, and
its variance deficit is
\begin{equation}
  \boxed{\;v_{k-1}-\Var(x_{k-1}) \;=\; \frac{s^2\sin^2(\Delta\theta_k)}{v_k}\;},
  \qquad \Delta\theta_k=\theta_k-\theta_{k-1}.
  \label{eq:deficit}
\end{equation}
\end{proposition}

\begin{proof}
Write $u=\cos\theta_k$, $w=\sin\theta_k$, $u'=\cos\theta_{k-1}$,
$w'=\sin\theta_{k-1}$, so $\abar_k=u^2$, $1-\abar_k=w^2$, $v_k=s^2u^2+w^2$. For a
Gaussian, $\E[\varepsilon\mid x_k]=w(x_k-um)/v_k$, hence
$\hat x_0=(x_k-w\hat\varepsilon)/u=(us^2x_k+w^2m)/v_k$. Substituting into
$x_{k-1}=u'\hat x_0+w'\hat\varepsilon$ gives $c_k=(u'us^2+w'w)/v_k$ and
$b_k=(u'w^2-w'wu)/v_k$. For the mean, $c_ku+b_k=u'(u^2s^2+w^2)/v_k=u'$, exactly.
For the variance, $\Var(x_{k-1})=c_k^2v_k=(c_kv_k)^2/v_k$ and
\begin{align}
  v_kv_{k-1}-(c_kv_k)^2
  &=(s^2u^2+w^2)(s^2u'^2+w'^2)-(s^2uu'+ww')^2\nonumber\\
  &=s^2\big[u^2w'^2+w^2u'^2-2uu'ww'\big]
  =s^2(uw'-wu')^2=s^2\sin^2(\Delta\theta_k),
\end{align}
so $v_{k-1}-c_k^2v_k=s^2\sin^2(\Delta\theta_k)/v_k$.
\end{proof}

\begin{corollary}[the criterion is about spread, not location]
\label{cor:pointmass}
The deficit is $O(s^2)$ and vanishes at $s=0$: DDIM is exact on a point mass
under \emph{every} schedule. Together with the exact mean transport, this says the
sampler's only Gaussian error is a loss of spread, and that questions about
schedules are questions about spread.
\end{corollary}

\begin{corollary}[what cosine equalizes]
\label{cor:cosines1}
If $s=1$ then $v_k\equiv1$ and \eqref{eq:deficit} reduces to
$\sin^2(\Delta\theta_k)$, which uniform $\theta$ makes equal at every step. So
the cosine schedule equalizes the per-step deficit \textbf{if and only if
$s=1$}. This is the geometric fact behind Figure~\ref{fig:coordinates}c, and it
is conditional in exactly the way Section~\ref{sec:cosine} anticipated.
\end{corollary}

\subsection{The cost is a metric on the $\leff$ axis, and it has finite length}

Composing \eqref{eq:deficit} along the chain, the total relative deficit is
\begin{equation}
  D \;=\; 1-\prod_{k}\Big(1-\frac{s^2\sin^2(\Delta\theta_k)}{v_kv_{k-1}}\Big)
    \;\approx\; \sum_k\frac{s^2\Delta\theta_k^2}{v_k^2}
  \qquad\text{for small steps.}
  \label{eq:total}
\end{equation}
The following reduction is the counterpart of Proposition~\ref{prop:clock} for
this second criterion.

\begin{proposition}[the cost is a sum of squared arc-length steps]
\label{prop:metric}
Let $g=\sig/s=e^{-\leff/2}$ be the noise level in units of the data scale. Then
\begin{equation}
  D \;\approx\; \sum_k\left(\frac{\Delta g_k}{1+g_k^2}\right)^{\!2}
    \;\approx\; \sum_k (\Delta\psi_k)^2,
  \qquad
  \psi=\arctan g=\arctan\frac{\sig}{s},
  \label{eq:sumsq}
\end{equation}
and in terms of effective log-SNR the arc-length element is
\begin{equation}
  \boxed{\;
  |\dd\psi| \;=\; \frac{\dd\leff}{4\cosh(\leff/2)}\;},
  \qquad
  \int_{-\infty}^{\infty}\frac{\dd\leff}{4\cosh(\leff/2)} \;=\; \frac{\pi}{2}.
  \label{eq:metric}
\end{equation}
So the cost depends on schedule and data only through $\leff$, it is a
\emph{metric} on the $\leff$ axis whose density peaks at $\leff=0$ and decays
like $e^{-|\leff|/2}$, and unlike the information metric of
Section~\ref{sec:info} its total length is \emph{finite}.
\end{proposition}

\begin{proof}
Substituting $\sig=\tan\theta$ gives $\dd\theta=\cos^2\!\theta\,\dd\sig$ and
$v=\cos^2\!\theta\,(s^2+\sig^2)$, so
$s^2\dd\theta^2/v^2=s^2\dd\sig^2/(s^2+\sig^2)^2$; putting $\sig=sg$ makes this
$\dd g^2/(1+g^2)^2$, which is $\dd\psi^2$ since $\psi=\arctan g$. For
\eqref{eq:metric}, $g=e^{-\leff/2}$ gives
$\dd g=-\tfrac12e^{-\leff/2}\dd\leff$ and
$1+g^2=1+e^{-\leff}$, so
$|\dd\psi|=\tfrac12 e^{-\leff/2}\dd\leff/(1+e^{-\leff})
=\dd\leff/[2(e^{\leff/2}+e^{-\leff/2})]$. The total is
$\int\mathrm{sech}(u/2)\,\dd u/4=\tfrac12\int\mathrm{sech}\,t\,\dd t=\pi/2$.
\end{proof}

Two things follow immediately, and they are the answers to ``why this notion of
uniformity?''

\begin{corollary}[the shift, again, from an independent premise]
\label{cor:forced2}
The transport cost is a functional of $\{\leff_k\}$ alone. Hence
Corollary~\ref{cor:forced} holds for it verbatim, and Recipe~\ref{rec:main} is
forced a second time --- now by the sampler's measured error rather than by
information flow. \emph{Two unrelated criteria pick out the same coordinate};
this, rather than any appeal to angles, is why $\leff$ is the right axis.
\end{corollary}

\begin{corollary}[equal steps, and the closed form]
\label{cor:equalsteps}
By \eqref{eq:sumsq} the design problem is: minimize $\sum_k(\Delta\psi_k)^2$
subject to $\sum_k\Delta\psi_k=\Psi$ fixed. By Cauchy--Schwarz (equivalently
convexity of $x\mapsto x^2$) the unique minimizer is $\Delta\psi_k=\Psi/T$, with
value
\begin{equation}
  D_{\min}=\frac{\Psi^2}{T}=\frac{\pi^2}{4T}
  \qquad\text{for the full path }\Psi=\pi/2,
  \label{eq:dmin}
\end{equation}
independent of $s$. \textbf{Equal steps in $\psi$ is not a desideratum but a
conclusion}: $\psi$ is by construction the arc length of the cost, so ``spend
the budget uniformly in $\psi$'' is a restatement of ``spread a fixed total cost
as evenly as possible over $T$ steps'', which is optimal because the cost is
quadratic in the step.
\end{corollary}

This settles the question of what makes one coordinate preferable to another.
Neither $\theta$ nor $\psi$ is privileged as an angle. What is privileged is the
\emph{cost density} \eqref{eq:metric}, which says: place steps in proportion to
where the sampler loses variance, namely near $\leff=0$, where the data variance
and the injected noise are comparable. Integrating that density produces a clock,
and that clock is $\psi$. Since the density is a function of $\leff$, so is the
clock; and since $\leff=\lambda+2\log s$, reading the clock requires knowing $s$.
At $s=1$ the clock is $\theta$, which is why cosine looked canonical.

\begin{figure}[t]
  \centering
  \includegraphics[width=0.86\textwidth]{figs/fig_geometry.pdf}
  \caption{\textbf{Picture of the conclusion, not the premise.} Plotting the
  marginal as (signal amplitude, noise amplitude)
  $=(s\sqrt{\abar},\sqrt{1-\abar})$ traces an ellipse of semi-axes $(s,1)$, on
  which $\psi$ of \eqref{eq:sumsq} is the polar angle and $\theta$ is only the
  eccentric angle. The optimum of Corollary~\ref{cor:equalsteps} spaces the polar
  angles evenly (right); the unshifted cosine schedule crowds them into the
  noise-dominated corner (left). When $s=1$ the ellipse is the unit circle and
  the two coincide.}
  \label{fig:geometry}
\end{figure}

\section{The recipe is exactly the reparametrization}
\label{sec:equiv}

We can now verify the first claim made in Section~\ref{sec:setup}.

\begin{proposition}[shift $=$ uniform $\psi$]
\label{prop:equiv}
Apply Recipe~\ref{rec:main} with estimate $\hat s$ to a target trajectory
$\lambda^{\rm target}_k$ whose angles are $\theta^{\rm target}_k
=\arctan\sig^{\rm target}_k$. If the true scale is $s=\hat s$, the resulting
schedule satisfies
\begin{equation}
  \psi_k \;=\; \theta^{\rm target}_k \qquad\text{for every }k,
\end{equation}
identically. In particular, shifting the cosine trajectory by $-2\log\hat s$
produces exactly the schedule that is uniform in $\psi$, which by
Corollary~\ref{cor:equalsteps} is optimal.
\end{proposition}

\begin{proof}
The recipe sets $\lambda_k=\lambda^{\rm target}_k-2\log\hat s$, i.e.\
$\sig_k^2=\hat s^2\,\sig^{{\rm target}\,2}_k$, i.e.\ $\sig_k=\hat s\,\sig^{\rm
target}_k$. Then
$\psi_k=\arctan(\sig_k/s)=\arctan(\sig^{\rm target}_k)=\theta^{\rm target}_k$
when $s=\hat s$. Uniform $\theta^{\rm target}$ therefore gives uniform $\psi$.
\end{proof}

So the two descriptions --- ``translate the log-SNR curve'' and ``make the polar
angle uniform'' --- are one statement, and the translation is the more primitive
of the two: it is what Corollaries~\ref{cor:forced} and \ref{cor:forced2} deliver
directly, whereas $\psi$ requires the extra step of integrating the cost
density. Numerically the identity holds to $10^{-14}$
(Section~\ref{sec:verify}).

\begin{remark}[what ``cannot be reproduced by reshaping'' means]
\label{rem:reshape}
Cosine's tunable knobs (the exponent, the offset $\delta$, a logit-SNR range)
change the map $k\mapsto\lambda_k$ while leaving the meaning of each $\lambda$
untouched. The correction of \eqref{eq:shift} instead relabels which physical
noise level a given $\lambda$ denotes. No amount of reshaping performs a
relabeling, which is the content of Chen's observation
\cite{chen2023importance}; it is a statement about the reshaping family, and
\emph{not} an obstruction. Recipe~\ref{rec:main} performs the relabeling in one
line.
\end{remark}

\begin{remark}[endpoints are a separate choice]
Proposition~\ref{prop:equiv} matches the \emph{shape} of the schedule exactly.
Where to start and stop --- the noisiest and cleanest levels actually visited ---
is an independent decision. Translating the trajectory as written carries
cosine's endpoints along with it, so the noisiest level moves from
$\sig_{T-1}$ to $\hat s\,\sig_{T-1}$. If instead one wants to keep the original
terminal SNR and re-space only the interior, replace $\Psi=\pi/2$ in
\eqref{eq:dmin} by $\Psi=\arctan(\sig_{T-1}/\hat s)$; the optimum is still equal
steps in $\psi$ and $D_{\min}=\Psi^2/T$.
\end{remark}

\begin{figure}[t]
  \centering
  \includegraphics[width=\textwidth]{figs/fig_shift.pdf}
  \caption{\textbf{The correction is a rigid translation.} (a) One nominal
  schedule induces a different effective clock \eqref{eq:shift} in each
  coordinate; the curves are vertical translates. The band marks
  $|\lambda|\le4$, where signal and noise are comparable and, by
  \eqref{eq:metric}, where the cost lives. (b) Recipe~\ref{rec:main} at $s=0.3$:
  the unshifted schedule's effective trajectory is $2\log(1/s)$ too noisy, and
  translating the nominal curve by the same amount restores it. (c) The
  translated schedule advances the cost-metric arc length $\psi$ at exactly
  constant rate (Proposition~\ref{prop:equiv}); the unshifted one does not.}
  \label{fig:shift}
\end{figure}

\section{The cost of the wrong scale}
\label{sec:robust}

One computation now yields both the price of ignoring $s$ and the price of
estimating it badly, because they are the same quantity.

\begin{proposition}[the $\cosh$ law]
\label{prop:cosh}
Build the schedule by Recipe~\ref{rec:main} from a full-path uniform-$\theta$
target using estimate $\hat s$, and run it on data of true scale $s$. Then, to
leading order in $1/T$,
\begin{equation}
  \boxed{\;
  D(\hat s) \;=\; \frac{\pi^2}{4T}\cdot\frac{1+\rho^2}{2\rho}
  \;=\; D_{\min}\cosh(\log\rho),
  \qquad \rho=\frac{s}{\hat s}. \;}
  \label{eq:cosh}
\end{equation}
\end{proposition}

\begin{proof}
The schedule has $\sig_k=\hat s\tan\phi_k$ with $\phi_k$ uniform,
$\Delta\phi=\pi/(2T)$, so $g_k=\sig_k/s=\tan\phi_k/\rho$. Then
$\Delta g_k=\rho^{-1}\sec^2\!\phi_k\,\Delta\phi$ and
$1+g_k^2=(\rho^2+\tan^2\phi_k)/\rho^2$, so by \eqref{eq:sumsq}
\begin{equation}
  D=\rho^2\Delta\phi^2\sum_k\frac{\sec^4\phi_k}{(\rho^2+\tan^2\phi_k)^2}
   \;\approx\;\rho^2\frac{\pi}{2T}\int_0^{\pi/2}\frac{\sec^4\phi\,\dd\phi}{(\rho^2+\tan^2\phi)^2}
   =\rho^2\frac{\pi}{2T}\int_0^{\infty}\frac{(1+t^2)\,\dd t}{(\rho^2+t^2)^2},
\end{equation}
using $t=\tan\phi$. Splitting $1+t^2=(\rho^2+t^2)+(1-\rho^2)$ and using
$\int_0^\infty(\rho^2+t^2)^{-1}\dd t=\pi/(2\rho)$ and
$\int_0^\infty(\rho^2+t^2)^{-2}\dd t=\pi/(4\rho^3)$ gives
$\pi(1+\rho^2)/(4\rho^3)$, hence $D=\pi^2(1+\rho^2)/(8T\rho)$.
\end{proof}

\begin{corollary}[the unshifted cosine schedule]
\label{cor:cosprice}
Taking $\hat s=1$ in \eqref{eq:cosh} --- that is, declining to shift at all ---
gives $\rho=s$ and
\begin{equation}
  D_{\cos}=\frac{\pi^2(1+s^2)}{8Ts}=D_{\min}\cosh(\log s).
  \label{eq:dcos}
\end{equation}
\end{corollary}

So there is no separate theory for ``the cosine schedule is suboptimal'' and
``a mis-estimated $\hat s$ is suboptimal'': the unshifted schedule is simply the
maximally mis-specified member of the family, the one that asserts $\hat s=1$.
The penalty $\cosh(\log\rho)$ is symmetric under $\rho\mapsto1/\rho$ and flat at
the optimum --- $\cosh(\log2)=1.25$, $\cosh(\log4)=2.125$ --- so a coarse
estimate captures nearly all of the available gain
(Figure~\ref{fig:penalty}). At $s=0.45$, for instance, not shifting costs
$1.34\times$, while shifting by an estimate wrong by a factor of two costs
$1.25\times$.

\begin{figure}[t]
  \centering
  \includegraphics[width=\textwidth]{figs/fig_penalty.pdf}
  \caption{\textbf{One law, two readings.} (a) Equation~\eqref{eq:cosh}: the
  penalty for using $\hat s$ when the truth is $s$, with measured composed
  deficits at $T=80$. (b) The same curve read as \eqref{eq:dcos}: the price of
  leaving the schedule unshifted is $\cosh(\log s)$, the case $\hat s=1$.}
  \label{fig:penalty}
\end{figure}

\begin{figure}[t]
  \centering
  \includegraphics[width=\textwidth]{figs/fig_allocation.pdf}
  \caption{\textbf{Where the budget goes.} (a) Noise level per step at $T=20$.
  (b) Fraction of steps with $\sig_k<s$, i.e.\ with $\leff_k>0$. By
  \eqref{eq:metric} the cost is symmetric about $\leff=0$, so the optimum spends
  exactly half the budget on each side --- which the shifted schedule does for
  every $s$, while the unshifted one starves the signal-dominated half as $s$
  shrinks.}
  \label{fig:allocation}
\end{figure}

\section{When one scale is not enough}
\label{sec:beyond}

Everything so far assumed a single number $s$. Three refinements matter, in
increasing order of practical importance.

\subsection{Anisotropic data: diagonal and general $\Sigma$}

For $\Sigma=\mathrm{diag}(s_1^2,\dots,s_d^2)$, equation \eqref{eq:shift} holds
per coordinate with its own offset $2\log s_i$, so a scalar schedule supplies one
translation where $d$ are wanted (Figure~\ref{fig:shift}a). For general
$\Sigma=U\Lambda U^{\!\top}$ nothing new happens, because the injected noise is
isotropic:
\begin{equation}
  V_k=\abar_kU\Lambda U^{\!\top}+(1-\abar_k)UU^{\!\top}
     =U\big[\abar_k\Lambda+(1-\abar_k)I\big]U^{\!\top},
  \label{eq:eig}
\end{equation}
so \textbf{the eigenbasis is the same at every noise level}, the forward process
never mixes principal axes, the exact DDIM map (built from $V_k^{-1}$ and
$\Sigma$) is diagonal in that basis, and the problem splits into $d$ scalar
problems with $s_i^2=\ell_i(\Sigma)$. Whitening by $\Sigma^{-1/2}$ is defined
only up to a rotation, and the choice is immaterial for the same reason.

Two responses, then. \emph{Standardize} --- diffuse
$z_0=\Sigma^{-1/2}(x_0-\mu)$ --- which sets every $s_i=1$, makes the target
trajectory optimal in every coordinate, and attains $d\cdot\pi^2/(4T)$; this is
the only exact fix, and it is what offline diffusion policies do with a
per-dimension normalizer \cite{chi2023diffusion}. Or \emph{compromise} on a
single $\hat s$, which is the next subsection with the expectation taken over
coordinates.

\subsection{Heterogeneous scale: the two scalars to track}
\label{sec:hetero}

In an RL setting the relevant width is not a single dataset statistic. What the
reverse chain must land on is the \emph{conditional} law $\pi(a\mid o)$ at the
current state, so the scale is a random variable $s_i(o)$ --- indexed by
coordinate and by state, and drifting as the policy trains. Averaging it and
plugging the average in is not obviously the right move. The $\cosh$ law settles
what to do, because it is exactly the right shape for taking expectations.

\begin{proposition}[best global shift under heterogeneity]
\label{prop:hetero}
Let $s$ denote the random variable $s_i(o)$ under the joint law of coordinate and
state-visitation, and let a single global $\hat s$ be used. Then by
\eqref{eq:cosh} the expected relative deficit is
\begin{equation}
  \frac{\mathcal L(\hat s)}{D_{\min}}
  \;=\; \E\big[\cosh(\log s-\log\hat s)\big]
  \;=\; \tfrac12\Big(\frac{M_+}{\hat s}+\hat s\,M_-\Big),
  \qquad
  M_+:=\E[s],\quad M_-:=\E[1/s],
  \label{eq:hetcost}
\end{equation}
which is convex in $\log\hat s$ and minimized at
\begin{equation}
  \boxed{\;
  \hat s_\star=\sqrt{\frac{M_+}{M_-}},
  \qquad
  \frac{\mathcal L(\hat s_\star)}{D_{\min}}=\sqrt{M_+M_-}\;\geq\;1,
  \;}
  \label{eq:hetopt}
\end{equation}
with equality if and only if $s$ is almost surely constant.
\end{proposition}

\begin{proof}
Expand $\cosh(b-\beta)=\tfrac12(e^{-\beta}e^{b}+e^{\beta}e^{-b})$ with
$b=\log s$, $\beta=\log\hat s$ and take expectations, giving \eqref{eq:hetcost}.
Differentiating in $\beta$ gives $-e^{-\beta}M_++e^{\beta}M_-=0$, hence
$\hat s_\star^2=M_+/M_-$ and the value $\sqrt{M_+M_-}$. Cauchy--Schwarz gives
$M_+M_-=\E[s]\,\E[1/s]\ge\big(\E[\sqrt s\cdot1/\sqrt s]\big)^2=1$ with equality
iff $s$ is degenerate.
\end{proof}

This answers the measurement question completely: \textbf{two scalars suffice,
and they are $\E[s]$ and $\E[1/s]$}, accumulated over states, coordinates and
whatever time window one trusts. Their ratio gives the shift to use; their
product gives the price of insisting on a single shift at all. It is worth
reading the second quantity carefully, because it is the useful one:

\begin{corollary}[$\sqrt{M_+M_-}$ is a decision rule]
\label{cor:decide}
$\sqrt{M_+M_-}$ is a floor: no choice of global $\hat s$ beats it. So if it is
close to $1$, tuning $\hat s$ is worthwhile and nothing else is; if it is large,
\emph{no} global scalar helps and the only remedy is to make the correction
state-dependent (Section~\ref{sec:perstate}).
\end{corollary}

\begin{remark}[the log-normal reading, and what to report]
\label{rem:lognormal}
If $\log s\sim\mathcal N(m,\tau^2)$ then $M_\pm=e^{\pm m+\tau^2/2}$, so
\begin{equation}
  \hat s_\star=e^{m}\ \ (\text{the geometric mean}),
  \qquad
  \frac{\mathcal L(\hat s_\star)}{D_{\min}}=e^{\tau^2/2}.
  \label{eq:lognormal}
\end{equation}
Equivalently, track $m=\E[\log s]$ and $\tau^2=\Var(\log s)$: the mean of the
logs sets the shift and their variance sets the irreducible penalty. This is the
precise sense in which one should ``average the logarithms, not the variances'':
using $\E[s]$ instead of $e^{m}$ costs an extra factor $\cosh(\tau^2/2)$, which is
$1.13$ at $\tau=1$ and $1.70$ at $\tau=1.5$. As a rule of thumb, $\tau\lesssim0.5$
(scale varying by less than about $1.6\times$) gives a floor below $1.13$ and a
global shift is essentially optimal; $\tau\gtrsim1$ gives a floor above $1.65$ and
calls for Section~\ref{sec:perstate}.
\end{remark}

\begin{figure}[t]
  \centering
  \includegraphics[width=\textwidth]{figs/fig_heterogeneity.pdf}
  \caption{\textbf{A heterogeneous scale.} (a) Equation~\eqref{eq:hetcost} for
  log-normal populations: the same $\cosh$ shape as
  Figure~\ref{fig:penalty}, lifted off $1$ by the floor $e^{\tau^2/2}$ (dotted),
  which no global $\hat s$ can reach below. (b) The floor against
  $\tau=\mathrm{sd}(\log s)$, with sampled values of $\sqrt{M_+M_-}$, and the
  extra cost of using the arithmetic mean $\E[s]$ in place of the geometric mean.}
  \label{fig:heterogeneity}
\end{figure}

\begin{remark}[measuring it in this code base]
\label{rem:measuring}
The quantity needed is the spread of the sampler's own output at fixed state.
Where $K$ parallel chains are already materialized as a
$[K,\text{batch},d]$ tensor, $s_i(o)^2$ is its variance over the $K$ axis, and
the two accumulators of \eqref{eq:hetopt} are batch means of $s_i(o)$ and
$1/s_i(o)$. Note this is not the same as the variance \emph{across} the $d$
action components at a single sample, which is what an ``action variance''
diagnostic typically reports; that statistic measures the shape of one action
vector, not the policy's conditional width, and is not the quantity here.
Since $K$ is finite, both accumulators need the correction of
Section~\ref{sec:finite}; Recipe~\ref{rec:estimate} is the complete procedure.
\end{remark}

\subsection{Finite samples: turning the cost into a loss}
\label{sec:finite}

Proposition~\ref{prop:hetero} averaged \eqref{eq:cosh} over a population of
states and coordinates, treating $s$ as \emph{known} at each. It never is: the
conditional width must itself be estimated, from the $K$ actions the sampler
draws at a given state. Section~\ref{sec:robust} priced a wrong $\hat s$ and
Proposition~\ref{prop:hetero} priced a single $\hat s$ for a spread-out
population; neither yet prices \emph{not knowing} $s$. Doing so is what converts
\eqref{eq:cosh} from a diagnostic into an objective one can minimize.

The observation that makes this easy is that $\cosh(\log s-\log\hat s)$ depends
on $(s,\hat s)$ only through the log-error, so \emph{it cannot tell where that
error came from}. A log-error of $0.3$ costs $1.045$ whether it arose because
states genuinely differ in width or because $\hat s$ was computed from five
samples at one state. Heterogeneity and estimation noise are quoted in the same
currency, and an accounting that charges for one but not the other is incomplete.

\begin{proposition}[the Bayes rule is the ratio of the $\pm1$ moments]
\label{prop:bayes}
Let $Q$ be any law expressing one's uncertainty about $s$: a population, a
sampling distribution, or a posterior. Then
\begin{equation}
  \E_Q\big[\cosh(\log s-\log\hat s)\big]
  \;=\;\tfrac12\Big(\frac{M_+}{\hat s}+\hat s\,M_-\Big),
  \qquad M_\pm=\E_Q\big[s^{\pm1}\big],
  \label{eq:bayesrisk}
\end{equation}
which is \eqref{eq:hetcost} verbatim. It is minimized at
$\hat s_\star=\sqrt{M_+/M_-}$ with value $\sqrt{M_+M_-}$, and at that point the
two terms of \eqref{eq:bayesrisk} are equal --- the rule balances the price of
over- against under-estimating. It is in general \emph{not} the mean, median or
mode of $Q$, and it exists only if both $M_+$ and $M_-$ are finite.
\end{proposition}

So the recipe for turning the cost into a loss is: \textbf{name the uncertainty,
take its $\pm1$ moments, form the ratio}. Nothing beyond the algebra of
Proposition~\ref{prop:hetero} is needed; what changes is only which $Q$ the
expectation is taken over. Two choices of $Q$ matter here, and they answer
different questions.

\subsubsection*{$Q$ as a posterior: the estimate to use at one state}

Let $a_1,\dots,a_K$ be iid draws of one action coordinate at a fixed state, and
$\hat s_K^2=\frac{1}{K-1}\sum_j(a_j-\bar a)^2$ the usual sample variance. If the
conditional is Gaussian then $\nu\hat s_K^2/s^2\sim\chi^2_\nu$ with
$\nu=K-1$ degrees of freedom.

\begin{proposition}[finite-$K$ Bayes estimate and its risk]
\label{prop:finiteK}
Under the scale-invariant prior $p(s)\propto1/s$ (flat in $\log s$, the natural
choice since the loss depends only on $\log s$), the posterior of $1/s^2$ is
$\mathrm{Gamma}(\nu/2,\ \nu\hat s_K^2/2)$, whence
\begin{equation}
  M_+=\hat s_K\sqrt{\tfrac\nu2}\,\frac{\Gamma(\frac{\nu-1}{2})}{\Gamma(\frac\nu2)},
  \qquad
  M_-=\frac{1}{\hat s_K}\sqrt{\tfrac2\nu}\,\frac{\Gamma(\frac{\nu+1}{2})}{\Gamma(\frac\nu2)},
\end{equation}
and therefore, using $\Gamma(\frac{\nu+1}{2})=\frac{\nu-1}{2}\Gamma(\frac{\nu-1}{2})$,
\begin{equation}
  \boxed{\;
  \hat s_\star=\hat s_K\sqrt{\frac{\nu}{\nu-1}},
  \qquad
  \frac{\mathcal L(\hat s_\star)}{D_{\min}}=\sqrt{f(\nu)},
  \quad
  f(\nu)=\frac{\Gamma(\frac{\nu-1}{2})\Gamma(\frac{\nu+1}{2})}{\Gamma(\frac{\nu}{2})^{2}}
  \;}
  \label{eq:finiteK}
\end{equation}
with $\sqrt{f(\nu)}=1+\frac{1}{4\nu}+O(\nu^{-2})$. The risk does not depend on the
observed $\hat s_K$: the factors of $\hat s_K$ cancel in $M_+M_-$.
\end{proposition}

Three things to read off \eqref{eq:finiteK}.

\begin{itemize}\itemsep2pt
  \item \textbf{The optimal estimate \emph{inflates} the sample width.} Not
  because the sample variance is biased --- $\hat s_K^2$ is unbiased --- but
  because the posterior of $\log s$ is right-skewed, and \eqref{eq:bayesrisk}
  charges symmetrically in $\log s$. Plugging in $\hat s_K$ under-shifts.
  \item \textbf{$K\ge3$ is required for this route.} At $\nu=1$,
  $\Gamma(\frac{\nu-1}{2})$ diverges: $M_+=\infty$, the risk is infinite, and no
  finite $\hat s$ is admissible. The obstruction is specific to the $-1$
  moment, not to small samples as such --- a per-state \emph{variance} is
  perfectly well behaved at $K=2$, and so is the log-space route of
  Proposition~\ref{prop:logroute} below, which is why that route is the one to
  prefer.
  \item \textbf{The penalty falls off fast.} $\sqrt{f(\nu)}$ is $1.25$ at $K=3$,
  $1.085$ at $K=5$, $1.036$ at $K=9$ and $1.017$ at $K=17$
  (Table~\ref{tab:finiteK}). Since not shifting at all typically costs
  $1.2$--$1.4$, $K$ of five to nine already recovers most of the available gain,
  and there is little reason to go further.
\end{itemize}

\subsubsection*{$Q$ as a sampling law: the global shift, debiased}

For the \emph{global} shift of Proposition~\ref{prop:hetero} one does not want a
per-state posterior; one wants to accumulate $M_\pm$ across states and
coordinates. The natural estimator replaces $s$ by $\hat s_K$ in both
accumulators, and this needs a correction, because
\begin{equation}
  \E[\hat s_K]=c(\nu)\,s,
  \quad
  \E[1/\hat s_K]=\frac{c'(\nu)}{s},
  \qquad
  c(\nu)=\sqrt{\tfrac2\nu}\frac{\Gamma(\frac{\nu+1}{2})}{\Gamma(\frac\nu2)},
  \quad
  c'(\nu)=\sqrt{\tfrac\nu2}\frac{\Gamma(\frac{\nu-1}{2})}{\Gamma(\frac\nu2)} .
  \label{eq:chibias}
\end{equation}
The two biases run in \emph{opposite} directions --- $c(\nu)<1$ while
$c'(\nu)>1$, since $\hat s_K$ understates $s$ on average while $1/\hat s_K$
overstates $1/s$ --- so they compound in the ratio and partially cancel in the
product. Since $c(\nu)c'(\nu)=f(\nu)$ and $c(\nu)/c'(\nu)=(\nu-1)/\nu$:

\begin{corollary}[what to correct, and by how much]
\label{cor:debias}
Write $\widehat{M}_\pm$ for the accumulators built from raw sample widths
$\hat s_K$. Then
\begin{equation}
  \sqrt{\frac{\widehat M_+}{\widehat M_-}}
  =\hat s_\star\sqrt{\frac{\nu-1}{\nu}},
  \qquad
  \sqrt{\widehat M_+\widehat M_-}
  =\sqrt{M_+M_-}\cdot\sqrt{f(\nu)} .
  \label{eq:debias}
\end{equation}
So the shift is recovered by multiplying by $\sqrt{\nu/(\nu-1)}$ --- the
\emph{same} factor as \eqref{eq:finiteK}, so one correction serves both readings
--- and the heterogeneity floor by dividing by $\sqrt{f(\nu)}$. Without the
second correction the measured floor is inflated by finite $K$ and
Corollary~\ref{cor:decide} will over-report the case for a state-dependent shift.
\end{corollary}

The distinction matters operationally. Estimation noise in a \emph{global} shift
averages out over states, so the $\sqrt{f(\nu)}$ term is a bias to be removed
from the reported floor, not a cost that is actually paid. In a \emph{per-state}
shift there is no averaging and $\sqrt{f(\nu)}$ is paid in full at every state.
That is the trade-off the next subsection has to resolve.

\subsubsection*{The route to actually use: accumulate in log space}

Proposition~\ref{prop:bayes} is what identifies the log-error as the thing that
costs; once that is known, the point estimate needs no further decision theory.
A global shift only ever wanted a location statistic for $\log s$, and there is a
simpler and better-behaved estimator of it than $\widehat M_\pm$.

\begin{proposition}[the log-space accumulators]
\label{prop:logroute}
With $\nu\hat s_K^2/s^2\sim\chi^2_\nu$,
\begin{equation}
  \boxed{\;
  \E\big[\log\hat s_K\big]=\log s+\beta(\nu),
  \qquad
  \beta(\nu)=\tfrac12\big(\psi(\nu/2)-\log(\nu/2)\big),
  \qquad
  \Var\big(\log\hat s_K\big)=\kappa(\nu),
  \;}
  \label{eq:logroute}
\end{equation}
both exactly, with $\psi$ the digamma function and $\kappa$ as in
\eqref{eq:kappa}. Hence tracking only the first two moments of
$\log\hat s_i(o)$ over states and coordinates gives
\begin{equation}
  \log\hat s_\star=\E\big[\log\hat s_K\big]-\beta(\nu),
  \qquad
  \tau^2=\Var\big(\log\hat s_K\big)-\kappa(\nu),
  \qquad
  \text{floor}=e^{\tau^2/2}.
  \label{eq:logaccum}
\end{equation}
\end{proposition}

Three advantages over $\widehat M_\pm$, and one caveat.

\begin{itemize}\itemsep2pt
  \item \textbf{Both corrections are exact additive constants} in the coordinate
  one is averaging in, and neither assumes anything about the population of $s$.
  The $\widehat M_\pm$ route instead needs two multiplicative corrections
  \eqref{eq:debias} applied to two separate accumulators.
  \item \textbf{One accumulator pair does both jobs.} The mean gives the shift
  and the variance gives $\tau^2$, hence the decision statistic of
  Corollary~\ref{cor:decide}, with no extra bookkeeping.
  \item \textbf{It is admissible at $K=2$.} $\beta(1)=-0.635$ and
  $\kappa(1)=1.234$ are finite, so two actions per state suffice for a global
  shift; the per-state noise is large ($\mathrm{sd}=1.11$ in $\log$) but it is
  unbiased and averages out over states. Nothing in
  Proposition~\ref{prop:finiteK} survives at $\nu=1$.
  \item \textbf{Caveat: it targets the geometric mean}, $e^{\E[\log s]}$, not
  $\sqrt{M_+/M_-}$. Writing $K_{\log s}$ for the cumulant generating function of
  $\log s$, $\log\sqrt{M_+/M_-}=\tfrac12[K_{\log s}(1)-K_{\log s}(-1)]$, so the
  two differ by the odd cumulants above the first:
  \begin{equation}
    \log\sqrt{M_+/M_-}-\E[\log s]
    =\frac{\kappa_3}{6}+\frac{\kappa_5}{120}+\cdots,
    \label{eq:oddcum}
  \end{equation}
  vanishing whenever $\log s$ is symmetric --- in particular for the log-normal
  case of Remark~\ref{rem:lognormal}, where the two coincide exactly. What that
  gap costs is settled exactly by Proposition~\ref{prop:selfsim} below.
\end{itemize}

The geometric mean is therefore \emph{not} the minimizer of the $\cosh$ cost
under heterogeneity, and it is worth being precise about what it is instead, and
about what the substitution costs. Both follow from one identity.

\begin{proposition}[the risk in $\log\hat s$ is again the $\cosh$ law]
\label{prop:selfsim}
For any law of $s$ with finite $M_\pm$, write $\beta=\log\hat s$ and
$\beta_\star=\log\sqrt{M_+/M_-}$. Then \eqref{eq:bayesrisk} is
\begin{equation}
  \boxed{\;
  \mathcal L(\beta)\;=\;\sqrt{M_+M_-}\,\cosh\!\big(\beta-\beta_\star\big)
  \;}
  \label{eq:selfsim}
\end{equation}
\emph{exactly}. Consequently the optimum is characterized by
\begin{equation}
  \E\big[\sinh(\log s-\beta_\star)\big]=0,
  \label{eq:sinhbalance}
\end{equation}
and \textbf{any} summary statistic $\beta$ costs exactly $\cosh(\beta-\beta_\star)$
relative to the optimum. In particular the geometric mean costs exactly
$\cosh(\Delta)$, with $\Delta=\beta_\star-\E[\log s]$ given by \eqref{eq:oddcum}.
\end{proposition}

\begin{proof}
Write $M_\pm=\sqrt{M_+M_-}\,e^{\pm\beta_\star}$, which is just the definition of
$\beta_\star$ together with $\sqrt{M_+M_-}$, and substitute into
$\tfrac12(e^{-\beta}M_++e^{\beta}M_-)$. Equation \eqref{eq:sinhbalance} is
$\mathcal L'(\beta_\star)=0$ written before the moments are collected.
\end{proof}

So Section~\ref{sec:robust}'s $\cosh$ law is self-similar: the penalty for a
misplaced \emph{summary of a population} has the same form as the penalty for a
misplaced $\hat s$ against a single $s$, because it is the same statement.

\begin{remark}[why the optimum is not the geometric mean]
\label{rem:whynotgeo}
Compare the two stationarity conditions: the geometric mean solves
$\E[\log s-\beta]=0$, i.e.\ it minimizes the \emph{squared} log-error, whereas
\eqref{eq:sinhbalance} weights deviations by $\sinh x=x+x^3/6+\cdots$. The
$\cosh$ cost therefore charges large log-errors far more than a quadratic does,
and its optimum is pulled toward the \emph{longer tail} of $\log s$ --- it buys
insurance against the states that would otherwise be badly mis-scheduled. Right-
skewed $\log s$ gives $\Delta>0$, i.e.\ a shift above the geometric mean.
Numerically, with $\Delta\approx\gamma_1\tau^3/6$ to leading order in the
log-skewness $\gamma_1$:
\begin{center}\small
\begin{tabular}{llll}
\toprule
$\tau$ & $\gamma_1$ & $\cosh\Delta$ (cost of geometric mean)
  & floor $e^{\tau^2/2}$ \\
\midrule
$0.6$ & $1$ & $1.0007$ & $1.197$ \\
$0.6$ & $2$ & $1.0043$ & $1.197$ \\
$1.0$ & $1$ & $1.020$  & $1.649$ \\
$1.0$ & $2$ & $1.328$  & $1.649$ \\
\bottomrule
\end{tabular}
\end{center}
So the substitution is harmless in the moderate regime but \emph{not}
unconditionally: at large $\tau$ and strong skew it becomes comparable to the
floor it is approximating. Note also that $\kappa_3/6$ is only the leading term
--- for a heavy right tail (log-skewness $2.8$) the true $\Delta=0.187$ against
$\kappa_3/6=0.103$.
\end{remark}

\begin{remark}[measuring $\Delta$ for free]
\label{rem:measuredelta}
One need not assume the approximation is safe. If $\widehat M_\pm$ are
accumulated alongside $\E[\log\hat s_K]$ --- three running scalars rather than
two, and available whenever $K\ge3$ --- then
$\Delta=\log\sqrt{\widehat M_+/\widehat M_-}+\tfrac12\log\frac{\nu}{\nu-1}
-\E[\log\hat s_K]+\beta(\nu)$ estimates the asymmetry of $\log s$ directly, and by
Proposition~\ref{prop:selfsim} $\cosh\Delta$ is exactly what the geometric mean
is costing. Report it; switch to $\sqrt{\widehat M_+/\widehat M_-}$ only if it is
materially above $1$. This keeps the log-space route as the default while making
its one approximation an observable rather than an assumption.
\end{remark}

\begin{remark}[two different corrections, two different questions]
\label{rem:twocorrections}
It is worth noting that $e^{-\beta(\nu)}\neq\sqrt{\nu/(\nu-1)}$: at $K=3$ they
are $1.335$ and $1.414$. They are not competing estimates of the same thing.
$\sqrt{\nu/(\nu-1)}$ is the Bayes hedge for \emph{one state's} $s$ under
\eqref{eq:bayesrisk}, and must pay for the posterior spread at that state;
$e^{-\beta(\nu)}$ merely makes $\log\hat s_K$ unbiased so that an average over
many states converges to $\E[\log s]$, the spread having been averaged away.
Using the former inside a global average over-shifts by
$\tfrac12\log\frac{\nu}{\nu-1}-\beta(\nu)$, a factor $1.06$ at $K=3$.
\end{remark}

\subsubsection*{Combining the two: shrinkage in log space}

Both effects are variances of $\log s$, so they simply add. Write
\begin{equation}
  \log\hat s_K(o)=\log s(o)+\epsilon,
  \qquad
  \Var(\epsilon)=\kappa(\nu):=\tfrac14\psi_1(\nu/2),
  \label{eq:kappa}
\end{equation}
$\psi_1$ being the trigamma function (unrelated to the polar angle $\psi$ of
\eqref{eq:sumsq}, which carries no subscript); $\kappa(\nu)\to1/(2\nu)$ for large
$\nu$, but the limit is poor where it matters --- at $K=3$ the exact
$\kappa=0.411$ against $1/(2\nu)=0.25$, understating the noise by $39\%$ --- so
use the closed form. With $\tau^2=\Var(\log s)$ as in
Remark~\ref{rem:lognormal}, the three candidate policies cost, in the log-normal
reading,
\begin{equation}
  \underbrace{e^{\tau^2/2}}_{\text{global}},
  \qquad
  \underbrace{e^{\kappa(\nu)/2}}_{\text{raw per-state}},
  \qquad
  \underbrace{e^{h/2}}_{\text{shrunk per-state}},
  \quad
  h=\Big(\tfrac{1}{\tau^2}+\tfrac{1}{\kappa(\nu)}\Big)^{-1},
  \label{eq:threecosts}
\end{equation}
the last obtained by shrinking toward the global geometric mean,
\begin{equation}
  \boxed{\;
  \log\hat s(o)=\log\hat s_\star+w\big(\log\hat s_K(o)-\log\hat s_\star\big),
  \qquad
  w=\frac{\tau^2}{\tau^2+\kappa(\nu)} .
  \;}
  \label{eq:shrink}
\end{equation}

\begin{corollary}[shrinkage is never the wrong choice]
\label{cor:shrink}
$h<\min(\tau^2,\kappa(\nu))$ strictly, so \eqref{eq:shrink} beats both the global
shift and the raw per-state estimate for every $\tau$ and every $K$. It
degenerates gracefully at both ends: $w\to0$ recovers the global shift when
samples are few or the population is homogeneous, and $w\to1$ recovers the raw
per-state estimate when $K$ is large. Raw per-state estimates beat a global shift
only when $\kappa(\nu)<\tau^2$, i.e.\ $\tau>\sqrt{\kappa(\nu)}$, which is
$0.64$ at $K=3$ and $0.27$ at $K=9$ --- so with few samples the ``more
adaptive'' choice is often the worse one, and \eqref{eq:shrink} is what removes
the need to guess.
\end{corollary}

\begin{table}[t]
  \centering
  \small
  \begin{tabular}{rrrrrr}
    \toprule
    $K$ & $\nu$ & $\sqrt{\nu/(\nu-1)}$ & $\sqrt{f(\nu)}$ & $\kappa(\nu)$
        & $\sqrt{\kappa(\nu)}$ \\
    & & \eqref{eq:finiteK} shift & \eqref{eq:finiteK} penalty & \eqref{eq:kappa}
        & break-even $\tau$ \\
    \midrule
    2  & 1  & $\infty$ & $\infty$ & --      & --     \\
    3  & 2  & 1.4142   & 1.2533   & 0.4112  & 0.6413 \\
    4  & 3  & 1.2247   & 1.1284   & 0.2337  & 0.4834 \\
    5  & 4  & 1.1547   & 1.0854   & 0.1612  & 0.4015 \\
    7  & 6  & 1.0954   & 1.0509   & 0.0987  & 0.3142 \\
    9  & 8  & 1.0690   & 1.0362   & 0.0710  & 0.2664 \\
    17 & 16 & 1.0328   & 1.0168   & 0.0333  & 0.1824 \\
    33 & 32 & 1.0160   & 1.0081   & 0.0161  & 0.1270 \\
    \bottomrule
  \end{tabular}
  \caption{Finite-sample constants. Column 3 multiplies the raw sample width (or
  the raw $\sqrt{\widehat M_+/\widehat M_-}$) to give the estimate to use;
  column 4 is the residual penalty and also the factor by which a measured
  heterogeneity floor is inflated; column 6 is the population spread above which
  a raw per-state shift beats a global one. $K=2$ is inadmissible:
  the $-1$ moment diverges.}
  \label{tab:finiteK}
\end{table}

\begin{recipe}[estimating $\hat s$ from $K$ samples per state]
\label{rec:estimate}
With the sampler's $K$ actions per state materialized as a $[K,\text{batch},d]$
tensor, and $\nu=K-1$:
\begin{enumerate}\itemsep1pt
  \item Per state and coordinate, $\hat s_i(o)^2=$ variance over the $K$ axis
  (\texttt{ddof}$=1$). Not over the $d$ axis (Remark~\ref{rem:measuring}).
  \item Accumulate the running \emph{mean and variance of
  $\log\hat s_i(o)$} over states and coordinates. Two scalars; nothing else is
  needed.
  \item Global shift and decision statistic by \eqref{eq:logaccum}:
  $\hat s_\star=\exp(\E[\log\hat s_K]-\beta(\nu))$ and
  $\tau^2=\Var(\log\hat s_K)-\kappa(\nu)$, both corrections being exact additive
  constants.
  \item Report the floor $e^{\tau^2/2}$. Below $\approx1.13$, stop and tune
  $\hat s_\star$; above $\approx1.65$, go per-state via \eqref{eq:perstate} with
  the shrinkage \eqref{eq:shrink}.
\end{enumerate}
Step 4 is the cheap experiment: it is a diagnostic computable at fixed $\hat s$,
before anything is allowed to affect training. If $K\ge3$, accumulate
$\widehat M_\pm$ as well and report $\cosh\Delta$ by
Remark~\ref{rem:measuredelta}: that is exactly what step 3's geometric mean
costs against the true optimum $\sqrt{M_+/M_-}$, and it is the only assumption in
the procedure worth checking rather than trusting. Switch step 3 to
$\sqrt{\widehat M_+/\widehat M_-}\cdot\sqrt{\nu/(\nu-1)}$ if it is materially
above $1$.
\end{recipe}

\begin{remark}[what is still idealized]
Two gaps between \eqref{eq:finiteK} and practice. The $\chi^2$ law assumes the
$K$ draws are iid Gaussian at fixed $o$; a diffusion sampler's outputs are
neither exactly Gaussian nor, if the chains share noise, exactly independent, so
$\nu$ is an upper bound on the effective degrees of freedom and the corrections
of Table~\ref{tab:finiteK} are best read as lower bounds. And $\tau^2$ estimated
by step 2 of Recipe~\ref{rec:estimate} conflates genuine state dependence with
anisotropy across the $d$ coordinates; the two call for different remedies
(\S\ref{sec:beyond} versus \S\ref{sec:perstate}), so it is worth accumulating
per-coordinate and pooled versions separately.
\end{remark}

\subsection{Making the correction state-dependent}
\label{sec:perstate}

When Corollary~\ref{cor:decide} says a global scalar is inadequate, note that
Proposition~\ref{prop:metric} leaves two ways to fix the ratio $\sig_k/s$: move
the numerator, or move the denominator. The schedule is usually a compile-time
constant, but the data is not. Standardizing per state,
\begin{equation}
  z_0=\frac{a-\hat\mu(o)}{\hat s(o)},
  \qquad\text{diffuse } z_0 \text{ under the fixed target schedule},
  \qquad a=\hat\mu(o)+\hat s(o)\,z_0,
  \label{eq:perstate}
\end{equation}
realizes a \emph{per-state} translation of $\leff$ while the schedule itself
never changes. This is the conditional version of the dataset-level normalizer of
\cite{chi2023diffusion}, and by Proposition~\ref{prop:equiv} it is exactly a
state-dependent shift. Its cost is an estimate of $\hat s(o)$ --- a small extra
head, or a running per-state statistic --- and its benefit is a portion of the
gap $\sqrt{M_+M_-}$ of Proposition~\ref{prop:hetero}: the whole gap only in the
limit of a perfectly known $\hat s(o)$, and in practice the fraction that
survives the estimation noise of \eqref{eq:threecosts}. Estimating $\hat s(o)$
from $K$ samples per state and using it raw can be \emph{worse} than a global
shift (Corollary~\ref{cor:shrink}); the shrinkage \eqref{eq:shrink} is what makes
this step safe.

\begin{remark}[why to sweep rather than trust, in online RL]
Two caveats specific to policy learning. First, the criterion of
Section~\ref{sec:transport} rewards fidelity to the target law, whereas in online
RL some of the sampler's variance loss is doing useful work as exploration noise
\cite{ren2024dppo}; sharpening the clean end of the schedule removes it. Second,
$s_i(o)$ is a property of the \emph{current} policy and moves during training, so
a frozen $\hat s$ is correct only for the window in which it was measured. Both
argue for treating $\hat s$ as a swept hyperparameter whose optimum is informed,
not determined, by \eqref{eq:hetopt}. The flatness of $\cosh$ is what makes this
tolerable.
\end{remark}

\subsection{Standardizing in the architecture: preconditioning and reuse under a drifting scale}
\label{sec:precond}

The two realizations above put the correction in the \emph{data}: standardize
the sample (\S\ref{sec:beyond}) or shift it per state (\S\ref{sec:perstate}).
Both bake $\hat s$ into the training target, so the previous remark's second
caveat --- that $s_i(o)$ drifts as the policy trains --- forces a moving target
and, in the naive reading, periodic retraining. A third realization removes the
offset \emph{inside the network} and makes tracking a drifting $\hat s$ exactly
free. It is the preconditioning of \cite{karras2022elucidating}, read through the
lens of \eqref{eq:shift0}.

Work in the signal frame $y_k:=x_k/\sqrt{\abar_k}=x_0+\sig_k\varepsilon$ on
centered data (the mean is removed by the exact coordinate shift of
\S\ref{sec:mean} and added back at the end). Write the denoiser as
\begin{equation}
  \hat x_0 \;=\; D(y,\sig)\;=\;c_{\rm skip}(\sig)\,y\;+\;c_{\rm out}(\sig)\,
  f_\theta\!\big(c_{\rm in}(\sig)\,y;\;\leff\big),
  \qquad
  c_{\rm in}=\tfrac{1}{\sqrt{s^2+\sig^2}},\;\;
  c_{\rm skip}=\tfrac{s^2}{s^2+\sig^2},\;\;
  c_{\rm out}=\tfrac{s\,\sig}{\sqrt{s^2+\sig^2}},
  \label{eq:precond}
\end{equation}
where --- and this is the part that matters here --- the network's noise argument
is the effective log-SNR $\leff=-2\log(\sig/s)$ of \eqref{eq:shift0}, not the bare
$\sig$ or the step index. The three coefficients are fixed by demanding that the
network input and its regression target each have unit variance at every $\sig$.

\begin{proposition}[scale invariance of the preconditioned target]
\label{prop:precond}
Write $x_0=s\,u$ with $u$ the standardized shape (zero mean, unit variance) and
$\rho=\sig/s$, so $\leff=-2\log\rho$. Under \eqref{eq:precond} both the network
input and its regression target are functions of $(u,\varepsilon,\rho)$ alone:
\begin{equation}
  c_{\rm in}\,y=\frac{u+\rho\,\varepsilon}{\sqrt{1+\rho^2}},
  \qquad
  \frac{x_0-c_{\rm skip}\,y}{c_{\rm out}}=\frac{\rho\,u-\varepsilon}{\sqrt{1+\rho^2}}.
  \label{eq:precond-inv}
\end{equation}
Hence the supervised problem $f_\theta$ solves --- map the former to the
conditional mean of the latter at each $\leff$ --- depends on the data only
through the law of $u$ and the value $\rho$, never on $s$. Its minimizer
$f_\theta^\star$ is invariant to the scale.
\end{proposition}

\begin{proof}
Substitute $x_0=su$ and $\sig=\rho s$ into \eqref{eq:precond}. Both denominators
are $\sqrt{s^2+\sig^2}=s\sqrt{1+\rho^2}$, and every numerator carries one power of
$s$, so $s$ cancels in each ratio, giving \eqref{eq:precond-inv}. The conditioning
$\leff=-2\log\rho$ carries no further $s$. The population risk at each $\leff$ is
thus a functional of the law of $(u,\varepsilon)$ only, and so is its minimizer.
\end{proof}

Equivalently, the optimal denoiser is homogeneous, $D(y;s)=s\,D_1(y/s)$ with $D_1$
the fixed $s{=}1$ denoiser: preconditioning \eqref{eq:precond} is exactly the
reparametrization that exposes this homogeneity to the optimizer, in the same way
Recipe~\ref{rec:main} exposes it to the sampler grid.

\begin{corollary}[reuse under a drifting scale]
\label{cor:reuse}
Let $\hat s$ (per coordinate, and $\hat\mu$ for the mean) be maintained online.
Updating it between steps recomputes only $c_{\rm in},c_{\rm skip},c_{\rm out}$
and the $\leff$ grid; by Proposition~\ref{prop:precond} the weights $\theta$
remain valid --- indeed remain the same optimum. There is no retraining and no
distribution shift for the network. Because the modeled clean law is
$D(\cdot;\hat s)=\hat s\,D_1(\cdot/\hat s)$, changing $\hat s$ moves no
probability mass the network must relearn: it is an exact reparametrization in
the sense of Recipe~\ref{rec:main}, so a policy defined through the sampler stays
valid across the update.
\end{corollary}

\begin{remark}[three ways to remove the same offset]
The $2\log s$ of \eqref{eq:shift0} can be taken out on the sampler grid
(Recipe~\ref{rec:main}: shift which $\lambda$'s are visited), in the data
(\S\ref{sec:beyond}, \S\ref{sec:perstate}: diffuse $(a-\hat\mu)/\hat s$), or in
the architecture (\eqref{eq:precond}: precondition and condition on $\leff$). The
three agree when $s$ is known and fixed. Preconditioning is the only one that is
at once per-$\sig$, per-coordinate --- replace $s$ by $\mathrm{diag}(s_i)$ and
$u=\Sigma^{-1/2}(x_0-\mu)$, everything staying diagonal by \eqref{eq:eig} --- and
free to track a drifting $\hat s$ (Corollary~\ref{cor:reuse}). Two boundaries are
worth stating. First, the invariance is to scale and mean \emph{only}: if the
standardized shape law of $u$ itself drifts, as it does while the policy improves,
that is genuine new signal the network must learn, preconditioning
notwithstanding. Second, it requires conditioning on $\leff$; a network keyed to
the step index or to bare $\sig$ forfeits Proposition~\ref{prop:precond} the
instant $s$ changes, because the same key then denotes a different $\rho$.
\end{remark}

\section{A different loss, a different summary: the advantage normalizer}
\label{sec:advnorm}

Proposition~\ref{prop:hetero} answered a population question: given one true
$s(o)$ per state and a budget of one global scalar, which scalar? The answer was
dictated entirely by the shape of the penalty, $\cosh$ of a log-ratio. A guided
diffusion policy carries a second global scalar of exactly this type --- the
divisor $\hat\sigma$ that normalizes the advantage before it tilts the policy ---
and it is worth working out, because the loss is different and so, it turns out,
is the summary. Nothing about ``average the per-state estimates in the log''
survives the change. What survives is the method: identify the coordinate in
which the cost is natural, then minimize.

\subsection{The tilt is an exponential family, so the cost is a Bregman divergence}

Let $\pi(a|o)$ be the current policy, $Q(o,a)$ the critic and $\beta$ the
guidance strength, and consider the tilted target
\begin{equation}
  p_{\hat\sigma}(a|o)\;\propto\;\pi(a|o)\,
  \exp\!\Big(\frac{\beta}{\hat\sigma}\,Q(o,a)\Big).
  \label{eq:tilt}
\end{equation}
Two observations fix the entire problem. First, \eqref{eq:tilt} is a natural
exponential family in $\eta=\beta/\hat\sigma$, with base measure $\pi(\cdot|o)$,
sufficient statistic $Q(o,\cdot)$ and log-partition
\begin{equation}
  \psi_o(\eta)=\log\E_{a\sim\pi(\cdot|o)}\big[e^{\eta Q(o,a)}\big],
  \label{eq:cgf}
\end{equation}
the cumulant generating function of $Q$ \emph{at fixed $o$}. Second, only that
conditional law can matter: replacing $Q(o,a)$ by $Q(o,a)+c(o)$ multiplies
numerator and denominator of \eqref{eq:tilt} by the same $e^{\beta c(o)/\hat\sigma}$
and leaves $p_{\hat\sigma}$ untouched. \textbf{The across-state spread of $Q$ is
pure gauge.} A normalizer built from the spread of $Q$ pooled over a batch of
states is therefore measuring the wrong quantity --- and by a factor
$\sqrt{1+\Sigma^2/\sigma_Q^2}$ that \emph{grows} as the critic's range $\Sigma$
across states grows, so it silently anneals the tilt away over training.

\begin{proposition}[cost of a wrong normalizer]
\label{prop:advkl}
Write $p_\eta$ for \eqref{eq:tilt}. For any $Q$ with finite CGF,
\begin{equation}
  \mathrm{KL}(p_{\eta_1}\|p_{\eta_2})
  =\psi_o(\eta_2)-\psi_o(\eta_1)-(\eta_2-\eta_1)\psi_o'(\eta_1),
  \label{eq:bregman}
\end{equation}
the Bregman divergence generated by the CGF, with $\psi_o'(\eta_1)=\E_{p_{\eta_1}}[Q]$.
If moreover $Q$ is Gaussian under $\pi(\cdot|o)$ with standard deviation
$\sigma_Q(o)$, then $\psi_o$ is quadratic, \eqref{eq:bregman} is exactly
quadratic and symmetric in its two arguments, and with $\eta_1=\beta/\hat\sigma$,
$\eta_2=\beta/\sigma_Q$,
\begin{equation}
  \boxed{\;
  \mathrm{KL}\big(p_{\hat\sigma}\,\|\,p_\star\big)
  =\frac{\beta^2}{2}\Big(\frac{\sigma_Q}{\hat\sigma}-1\Big)^{2}
  =\frac{\beta^2}{2}\big(e^{u}-1\big)^{2},\qquad
  u=\log\frac{\sigma_Q}{\hat\sigma}.
  \;}
  \label{eq:advcost}
\end{equation}
\end{proposition}

\begin{proof}
$\log(p_{\eta_1}/p_{\eta_2})=(\eta_1-\eta_2)Q-\psi_o(\eta_1)+\psi_o(\eta_2)$;
take $\E_{p_{\eta_1}}$ and use $\psi_o'(\eta_1)=\E_{p_{\eta_1}}[Q]$. For Gaussian
$Q$, $\psi_o(\eta)=\eta\mu_Q+\tfrac12\eta^2\sigma_Q^2$, whose Bregman divergence
is $\tfrac12\sigma_Q^2(\eta_2-\eta_1)^2$; substituting the two $\eta$'s cancels
$\sigma_Q^2$ against $\beta^2(1/\hat\sigma-1/\sigma_Q)^2$.
\end{proof}

\begin{remark}[why $\sigma_Q$ is the target: it equalizes the KL spend]
For Gaussian $Q$ the tilt is a pure mean shift, and the ideally normalized
target satisfies $\mathrm{KL}(p_\star\|\pi)=\beta^2/2$ \emph{at every state,
independent of $\sigma_Q$}. That is what advantage normalization is for:
dividing by the within-state $\sigma_Q(o)$ is exactly the choice that turns
$\beta$ into a state-independent KL budget. Equation~\eqref{eq:advcost} then says
that a wrong divisor spends the wrong budget at that state.
\end{remark}

\begin{remark}[the geometry is \emph{not} the $\cosh$ geometry]
\label{rem:advasym}
$(e^u-1)^2$ is not symmetric in $u$, and the asymmetry is the operative fact. As
$\hat\sigma\to\infty$ we get $\eta\to0$, $p_{\hat\sigma}\to\pi$, and the cost
saturates at
\begin{equation}
  \sup_{u<0}\ \frac{\beta^2}{2}\big(e^u-1\big)^2=\frac{\beta^2}{2}
  =\mathrm{KL}(\pi\|p_\star),
\end{equation}
so \emph{under}-tilting can do no worse than abandoning the tilt altogether.
Over-tilting is unbounded, growing like $e^{2u}$: at $u=\pm1$ the two costs are
$2.95$ and $0.40$, at $u=\pm2$ they are $40.8$ and $0.75$. Where a wrong $\hat s$
is penalized evenhandedly by $\cosh$, a wrong $\hat\sigma$ must be erred on the
\emph{high} side.
\end{remark}

\subsection{The population minimizer is the contraharmonic mean}

\begin{proposition}
\label{prop:contraharm}
Let $\sigma_Q(o)$ vary across states with $\E[\sigma_Q^2]<\infty$. The single
global divisor minimizing $\E_o[\mathrm{KL}]$ under \eqref{eq:advcost} is the
\emph{contraharmonic mean}
\begin{equation}
  \boxed{\;\hat\sigma_\star=\frac{\E[\sigma_Q^2]}{\E[\sigma_Q]}
  =\E[\sigma_Q]\big(1+\mathrm{CV}^2\big)\;}
  \qquad
  \mathrm{CV}=\frac{\mathrm{sd}(\sigma_Q)}{\E[\sigma_Q]},
  \label{eq:contraharm}
\end{equation}
and the residual risk at that choice is
\begin{equation}
  \E_o[\mathrm{KL}]_\star=\frac{\beta^2}{2}\cdot\frac{\mathrm{CV}^2}{1+\mathrm{CV}^2}.
  \label{eq:advfloor}
\end{equation}
\end{proposition}

\begin{proof}
Substitute $x=1/\hat\sigma$. The risk is
$\tfrac{\beta^2}{2}\big(\E[\sigma_Q^2]x^2-2\E[\sigma_Q]x+1\big)$, a convex
parabola in $x$ minimized at $x_\star=\E[\sigma_Q]/\E[\sigma_Q^2]$, which is
\eqref{eq:contraharm}. Its value is
$\tfrac{\beta^2}{2}\big(1-\E[\sigma_Q]^2/\E[\sigma_Q^2]\big)$, and
$\E[\sigma_Q]^2/\E[\sigma_Q^2]=1/(1+\mathrm{CV}^2)$ by definition of $\mathrm{CV}$;
the second form of \eqref{eq:contraharm} is the same identity.
\end{proof}

Unlike the discussion following Proposition~\ref{prop:selfsim}, the aggregation
is not in question here. KL is additive over states, so $\E_o$ under the state
distribution's own weights simply \emph{is} the KL between the joint policies;
there is no equal-weight-versus-$s^2$-weight choice to make.

\begin{remark}[three means, and which of them over-tilt]
\label{rem:threemeans}
By the power-mean inequality the arithmetic, quadratic and contraharmonic means
are ordered,
\begin{equation}
  \E[\sigma_Q]\;\le\;\sqrt{\E[\sigma_Q^2]}\;\le\;\frac{\E[\sigma_Q^2]}{\E[\sigma_Q]},
  \qquad\text{in the ratios } 1:\sqrt{1+\mathrm{CV}^2}:1+\mathrm{CV}^2 .
\end{equation}
Both of the first two are therefore \emph{too small}: they over-tilt, which by
Remark~\ref{rem:advasym} is the expensive direction. This matters in practice
because $\sqrt{\E_o \Var_K(Q)}$ --- the natural thing to write down from a batch
of per-state sample variances --- is the quadratic mean, short of the optimum by
exactly $\sqrt{1+\mathrm{CV}^2}$. Excess over the minimum \eqref{eq:advfloor}:
\begin{center}\small
\begin{tabular}{lccc}
\toprule
$\mathrm{CV}$ & $\E_o[\mathrm{KL}]_\star/(\beta^2/2)$
  & excess of $\sqrt{\E[\sigma_Q^2]}$ & excess of $\E[\sigma_Q]$ \\
\midrule
$0.2$ & $0.038$ & $+1.0\%$  & $+4.0\%$ \\
$0.5$ & $0.200$ & $+5.6\%$  & $+25.0\%$ \\
$1.0$ & $0.500$ & $+17.1\%$ & $+99.9\%$ \\
\bottomrule
\end{tabular}
\end{center}
\end{remark}

\begin{corollary}[what to accumulate]
\label{cor:advaccum}
From $K$ draws of $Q$ at a fixed state let $\hat\sigma_K(o)$ be the sample
standard deviation, $\nu=K-1$. Then $\E[\hat\sigma_K^2]=\sigma_Q^2$ by Bessel and
$\E[\hat\sigma_K]=c_4(\nu)\sigma_Q$ with
$c_4(\nu)=\sqrt{2/\nu}\,\Gamma(\tfrac{\nu+1}{2})/\Gamma(\tfrac\nu2)$, so running
averages of $\hat\sigma_K$ and $\hat\sigma_K^2$ over states give
\begin{equation}
  \hat\sigma_\star=c_4(\nu)\,\frac{\overline{\hat\sigma_K^2}}{\overline{\hat\sigma_K}},
  \qquad
  \mathrm{CV}^2=c_4(\nu)^2\,
  \frac{\overline{\hat\sigma_K^2}}{\overline{\hat\sigma_K}^{\,2}}-1 .
  \label{eq:advaccum}
\end{equation}
Two accumulators in \emph{linear} space, both corrections exact multiplicative
constants. Unlike the $M_-$ route of Corollary~\ref{cor:debias} nothing diverges
at $K=2$, where $c_4(1)=\sqrt{2/\pi}=0.7979$.
\end{corollary}

As in \eqref{eq:logaccum}, unbiasedness is the only criterion that matters for
either accumulator, and for the same reason: a running average over states kills
estimator variance but passes estimator bias through untouched. Ignoring $c_4$
at $K=2$ inflates $\hat\sigma$ by $1/c_4=1.25$, a $25\%$ under-tilt --- the cheap
direction by Remark~\ref{rem:advasym}, but free to fix.

\begin{remark}[the two scalars, side by side]
\label{rem:twoscalars}
The contrast is the point of this section: the same population question, asked
under two different losses, gives two different summaries and two different
accumulators.
\begin{center}\small
\begin{tabular}{lll}
\toprule
 & schedule shift $\hat s$ & advantage divisor $\hat\sigma$ \\
\midrule
per-state target & policy width $s(o)$ & within-state $\sigma_Q(o)$ \\
cost of an error & $\cosh u$, \eqref{eq:cosh}
  & $\tfrac{\beta^2}{2}(e^u-1)^2$, \eqref{eq:advcost} \\
symmetric in $u$? & yes & no --- err high \\
global optimum & $\sqrt{M_+/M_-}\approx$ geometric mean
  & contraharmonic $\E[\sigma^2]/\E[\sigma]$ \\
accumulate & $\log\hat s_K,\ (\log\hat s_K)^2$
  & $\hat\sigma_K,\ \hat\sigma_K^2$ \\
in which space & log & linear \\
debiased by & $-\beta(\nu)$, additive & $\times c_4(\nu)$, multiplicative \\
dof from $K$ & $d(K-1)$, pooling coordinates & $K-1$ \\
heterogeneity floor & $e^{\tau^2/2}$
  & $\tfrac{\beta^2}{2}\mathrm{CV}^2/(1+\mathrm{CV}^2)$ \\
\bottomrule
\end{tabular}
\end{center}
Both floors are diagnostics computable at a fixed scalar, before either estimate
is allowed to drive anything --- which is how they should first be run.
\end{remark}

\section{The mean plays no role}
\label{sec:mean}

\begin{proposition}
None of $V_k$, $\leff$, the information rate \eqref{eq:clockdep}, the deficit
\eqref{eq:deficit}, the cost metric \eqref{eq:metric}, the optimum
\eqref{eq:dmin} or the penalty \eqref{eq:cosh} depends on $\mu$.
\end{proposition}

The algebra appeared three times already: $V_k=\abar_k\Sigma+(1-\abar_k)I$ in
\eqref{eq:marginal} contains no $\mu$; $\mu$ dropped out of
Proposition~\ref{prop:clock} because a known additive constant carries no
information; and $c_k$ in Proposition~\ref{prop:deficit} is built from $\abar$,
$s$ and $V_k$ alone, with DDIM transporting the mean exactly. Three
complementary reasons this had to happen:

\paragraph{1. A translation is free and invertible.} Replacing $x_0$ by
$x_0+\mu$ replaces $x_k$ by $x_k+\sqrt{\abar_k}\mu$, a known deterministic
bijection at every level. It changes no entropy, no mutual information, and no KL
between centered marginals. A schedule allocates steps to \emph{removing
uncertainty}, and a mean carries none.

\paragraph{2. SNR is a ratio of variances.} The ``signal'' in a denoising
problem is the part of $x_0$ that is not predictable in advance. A mean is
predictable by definition and a network absorbs it into an output bias, so
$\snr=\Var(\text{signal})/\Var(\text{noise})$ correctly excludes it.

\paragraph{3. The recursions are decoupled.} Mean and covariance propagate
through \eqref{eq:forward} separately, $\E[x_k]=\sqrt{\abar_k}\E[x_0]$ and
$V_k=\abar_kV_0+(1-\abar_k)I$, with no cross term.

\paragraph{Two caveats.} Defining SNR from the raw second moment
$\E\|x_0\|^2=\|\mu\|^2+\Tr\Sigma$ does make it move with $\mu$, but that is the
wrong quantity for scheduling, by reason 2; the $\sigma_{\rm data}$ of EDM
\cite{karras2022elucidating} is a centered standard deviation. And the mean
\emph{does} enter a different notion of burden --- for scalar
$p_0=\mathcal N(m,s^2)$,
\begin{equation}
  2\,\mathrm{KL}(q_{k-1}\|q_k)\approx\Delta\theta_k^2
  \Big[\underbrace{\tfrac{\sin^2(2\theta)(1-s^2)^2}{2v^2}}_{\text{variance motion}}
  +\underbrace{\tfrac{m^2\sin^2\theta}{v}}_{\text{mean motion}}\Big],
  \label{eq:kl}
\end{equation}
so equalizing \emph{how far the marginal moves} per step is $\mu$-dependent,
while equalizing the transport deficit is not. The two disagree: the variance
term in \eqref{eq:kl} peaks at $\snr=1/s^2$ and vanishes at both ends, whereas the
cost density \eqref{eq:metric} peaks at $\snr^{\rm eff}=1$. Centering the data
remains good practice --- a large $\|\mu\|$ forces the network to emit a large
offset --- but it is a conditioning question, not a scheduling one.

\section{Numerical verification}
\label{sec:verify}

Closed forms were checked against direct computation using the schedules built by
the code base, in \texttt{scripts/equal\_burden\_}\-\texttt{schedule.py}, with
figures from \texttt{notes/snr\_schedule\_shift/}\-\texttt{make\_figures.py}. The
finite-sample results of Section~\ref{sec:finite} --- and the constants of
Table~\ref{tab:finiteK} --- are checked by Monte Carlo against their closed forms
in \texttt{notes/snr\_schedule\_shift/}\-\texttt{verify\_finite\_k.py}.

\begin{table}[h]
  \centering
  \small
  \begin{tabular}{llll}
    \toprule
    Claim & Prediction & Measured & Agreement \\
    \midrule
    I-MMSE, Gaussian \eqref{eq:clockdep}
      & $\dd I/\dd\leff=\tfrac12\sigma(\leff)$ & finite differences
      & $4\times10^{-11}$ \\
    I-MMSE, two-point prior
      & $\int\!\dd I=\log2=0.693147$ & $0.693147$ & $<10^{-6}$ \\
    Per-step deficit \eqref{eq:deficit}
      & $s^2\sin^2(\Delta\theta_k)/v_k$ & direct affine map
      & $4.4\times10^{-16}$ \\
    Shift $=$ uniform $\psi$ (Prop.~\ref{prop:equiv})
      & $\psi_k-\theta^{\rm target}_k=0$ & $1.1\times10^{-14}$ & exact \\
    $D_{\min}$ \eqref{eq:dmin}, $T=300$
      & $T\!\cdot\!D=\pi^2/4=2.4674$ & $2.4735$ & $0.2\%$ \\
    \ \ same, $s=0.1,0.2,0.45,0.8,1.0$
      & identical for all $s$ & $2.4734$--$2.4736$ & $10^{-4}$ \\
    $\cosh$ law \eqref{eq:cosh}, $\rho=2$
      & $1.250$ & $1.243$ & $0.6\%$ \\
    $\cosh$ law \eqref{eq:cosh}, $\rho=4$
      & $2.125$ & $2.080$ & $2.1\%$ \\
    Unshifted cosine \eqref{eq:dcos}, $s=0.45$
      & $1.336$ & $1.386$ & $3.7\%$ \\
    Heterogeneity floor \eqref{eq:hetopt}
      & $\sqrt{M_+M_-}$ & grid search over $\hat s$
      & $<10^{-3}$ \\
    \ \ log-normal form \eqref{eq:lognormal}
      & $e^{\tau^2/2}$, $\hat s_\star=e^m$ & sampled populations & $<10^{-3}$ \\
    Finite-$K$ shift \eqref{eq:finiteK}, $K=3$
      & $\sqrt{\nu/(\nu-1)}=1.41421$ & posterior risk minimizer $1.41553$
      & $0.1\%$ \\
    \ \ residual risk, $K=3$--$33$
      & $\sqrt{f(\nu)}$ & $\E_Q\cosh$ at the optimum & $<3\times10^{-4}$ \\
    Plug-in bias \eqref{eq:debias}, $K=3$
      & $\sqrt{(\nu-1)/\nu}=0.70711$ & $0.70692$ & $0.03\%$ \\
    \ \ debiased floor, $K=3,5,9,17$
      & $e^{\tau^2/2}=1.1972$ & $1.1964$--$1.1968$ & $0.06\%$ \\
    Estimation noise \eqref{eq:kappa}
      & $\kappa(\nu)=\tfrac14\psi_1(\nu/2)$ & $\Var(\log\hat s_K)$ sampled
      & $<0.1\%$ \\
    Shrinkage weight \eqref{eq:shrink}
      & $w=\tau^2/(\tau^2+\kappa)$ & MC risk minimizer
      & $\le0.3\%$ in risk \\
    Log-bias \eqref{eq:logroute}, $K=2$--$33$
      & $\beta(\nu)=\tfrac12(\psi(\tfrac\nu2)-\log\tfrac\nu2)$
      & $\E[\log\hat s_K]$ sampled & $<10^{-3}$ \\
    \ \ global shift \eqref{eq:logaccum}, $K=2$
      & $e^{\E[\log s]}=0.39976$ & $0.39964$ & $0.03\%$ \\
    Odd-cumulant gap \eqref{eq:oddcum}
      & $\kappa_3/6=\pm0.0494$ & $\pm0.0556$ & $\kappa_5/120$ residual \\
    \ \ its cost, $|\mathrm{skew}(\log s)|=1.4$
      & $\cosh\Delta=1.00155$ & $1.00155$ & $<10^{-6}$ \\
    Self-similar risk \eqref{eq:selfsim}
      & $\mathcal L(\beta)=\sqrt{M_+M_-}\cosh(\beta-\beta_\star)$
      & 4 populations $\times$ 41 $\beta$ & $3\times10^{-15}$ \\
    \ \ sinh balance \eqref{eq:sinhbalance}
      & $\E[\sinh(\log s-\beta_\star)]=0$ & sampled & $<10^{-15}$ \\
    \addlinespace
    Bregman form \eqref{eq:bregman}
      & $\psi(\eta_2)-\psi(\eta_1)-\Delta\eta\,\psi'(\eta_1)$
      & direct KL, non-Gaussian $Q$ & $2\times10^{-16}$ \\
    \ \ Gaussian cost \eqref{eq:advcost}
      & $\tfrac{\beta^2}{2}(\sigma_Q/\hat\sigma-1)^2$ & exact KL of the two tilts
      & $0$ \\
    Contraharmonic \eqref{eq:contraharm}
      & $\E[\sigma^2]/\E[\sigma]$ & MC risk minimizer & $<10^{-4}$ \\
    \ \ floor \eqref{eq:advfloor}
      & $\mathrm{CV}^2/(1+\mathrm{CV}^2)$ & MC risk at optimum & $<10^{-3}$ \\
    \ \ accumulators \eqref{eq:advaccum}, $K=2$
      & $c_4\overline{\hat\sigma^2}/\overline{\hat\sigma}=1.358$ & true $1.361$ & $0.2\%$ \\
    \bottomrule
  \end{tabular}
  \caption{Predictions against direct numerical evaluation; the last five rows are
  produced by \texttt{verify\_adv\_norm.py}, the rest by
  \texttt{verify\_finite\_k.py}. Residual gaps in the
  $D$ rows are the $O(\Delta\psi^2)$ error of the small-step step of
  Proposition~\ref{prop:metric}, which needs $T\gg\pi/(2s)$; the two exact
  identities hold to machine precision.}
\end{table}

\begin{figure}[h]
  \centering
  \includegraphics[width=\textwidth]{figs/fig_verification.pdf}
  \caption{\textbf{Convergence of the closed forms.} (a) Unshifted, $T\!\cdot\!D$
  tends to $\pi^2(1+s^2)/(8s)$ (dotted), a different constant for each $s$. (b)
  With Recipe~\ref{rec:main} it tends to $\pi^2/4$ for every $s$; the four curves
  coincide.}
  \label{fig:verification}
\end{figure}

\section{Summary}

\begin{enumerate}\itemsep3pt
  \item \textbf{The clock is exact, not asymptotic.} By I-MMSE,
  $\dd I/\dd\lambda=\tfrac12\gamma\,\mmse(\gamma)$ at every noise level for every
  clean law. The two factors are ``how much is still unknown'' and ``how good the
  channel is''; they trade off, and their product is the rate. For Gaussian data
  it is exactly $\tfrac12\sigma(\leff)$, tending to $\tfrac12$ at high SNR --- the
  usual high-SNR statement is the plateau of an exact curve.
  \item \textbf{The rate is a function of $\leff=\lambda+2\log s$ alone}
  (Proposition~\ref{prop:clock}), for any data shape. Since a schedule controls
  $\lambda$, it is under-determined by one constant until $s$ is named; naming
  nothing asserts $s=1$. This, not any statement about angles, is the primal
  desideratum.
  \item \textbf{Equal information per step cannot be the whole criterion}, because
  for continuous data the total information is infinite and accumulates at
  constant rate over an unbounded stretch of $\leff$. It fixes the coordinate, not
  the spacing. Cosine responds by compactifying (uniform $\theta$, total range
  $\pi/2$) rather than truncating, and coincides with uniform $\lambda$ only near
  $\snr=1$, where $|\dd\lambda/\dd\theta|=4$ is minimal.
  \item \textbf{The sampler's transport cost supplies the spacing.} Its exact
  per-step value is $s^2\sin^2(\Delta\theta_k)/v_k$, it also depends only on
  $\leff$ --- so it forces the same shift from an independent premise --- and it is
  a metric $\dd\leff/[4\cosh(\leff/2)]$ on the $\leff$ axis of \emph{finite} total
  length $\pi/2$, concentrated where signal and noise are comparable.
  \item \textbf{Uniformity in $\psi$ is a conclusion, not an assumption.} Because
  the cost is $\sum_k(\Delta\psi_k)^2$ with $\psi$ the arc length of that metric,
  minimizing at fixed $T$ is just ``spread a fixed total evenly'', giving
  $D_{\min}=\pi^2/(4T)$ independent of $s$. At $s=1$ the arc length is $\theta$,
  which is why cosine looked canonical.
  \item \textbf{The fix is Recipe~\ref{rec:main}}: $\lambda_k=\lambda^{\rm
  target}_k-2\log\hat s$. It is exactly the uniform-$\psi$ reparametrization
  (Proposition~\ref{prop:equiv}, verified to $10^{-14}$), it is exactly equivalent
  to standardizing the data, and it is a relabeling rather than a reshaping ---
  which is why cosine's shape knobs cannot express it, and why implementing it
  directly is trivial.
  \item \textbf{One law covers all the error.} $D=D_{\min}\cosh\log(s/\hat s)$;
  the unshifted schedule is the case $\hat s=1$, costing $\cosh(\log s)$. The
  penalty is flat --- a factor-of-two error costs $1.25\times$ --- so a coarse
  frozen estimate captures nearly all of the gain.
  \item \textbf{Anisotropy} reduces to $d$ scalar problems in the eigenbasis of
  $\Sigma$, which is noise-level independent because the injected noise is
  isotropic. Standardizing is the only exact fix; otherwise use the scalar
  compromise below.
  \item \textbf{For heterogeneous or state-dependent $s$, track two scalars:}
  $M_+=\E[s]$ and $M_-=\E[1/s]$ over states and coordinates. Then
  $\hat s_\star=\sqrt{M_+/M_-}$ and the irreducible penalty is
  $\sqrt{M_+M_-}\ge1$ --- equivalently report $m=\E[\log s]$ and
  $\tau^2=\Var(\log s)$, giving $\hat s_\star=e^m$ and floor $e^{\tau^2/2}$. The
  floor is a decision rule: small means tune $\hat s$, large means no global
  scalar helps and the correction must be made per state, which
  \eqref{eq:perstate} does by rescaling the data while leaving the schedule fixed.
  \item \textbf{Finite samples cost the same currency as heterogeneity.} Because
  the penalty sees only the log-error, it cannot distinguish genuine spread from
  estimation noise, so the $\pm1$-moment rule applies verbatim with $Q$ a
  posterior (Proposition~\ref{prop:bayes}). From $K$ draws per state it gives
  $\hat s_\star=\hat s_K\sqrt{\nu/(\nu-1)}$, $\nu=K-1$ --- an \emph{inflation} of
  the sample width, since $\log s$ has a right-skewed posterior --- at residual
  cost $\sqrt{f(\nu)}\approx1+1/(4\nu)$; that route needs $K\ge3$, since the $-1$
  moment diverges at $\nu=1$. \textbf{For a global shift, however, no loss is
  needed once the coordinate is known}: average $\log\hat s_K$ and subtract the
  exact bias $\beta(\nu)=\tfrac12(\psi(\nu/2)-\log(\nu/2))$
  (Proposition~\ref{prop:logroute}). Two accumulators, both corrections additive,
  admissible at $K=2$, and the variance gives $\tau^2$ for free. Its one
  approximation is that it targets the geometric mean, which is the minimizer of
  the \emph{squared} log-error rather than of $\cosh$: the true optimum solves
  $\E[\sinh(\log s-\beta_\star)]=0$ and sits toward the longer tail of $\log s$.
  Because the risk in $\log\hat s$ is \emph{itself} exactly a $\cosh$
  (Proposition~\ref{prop:selfsim}), the substitution costs exactly $\cosh\Delta$
  with $\Delta$ the odd-cumulant sum \eqref{eq:oddcum} --- negligible for
  moderate skew, comparable to the floor for strong skew, and cheap to measure
  rather than assume (Remark~\ref{rem:measuredelta}). A per-state shift should
  shrink toward the global one with weight $\tau^2/(\tau^2+\kappa(\nu))$, which is
  never worse than either extreme.
  \item \textbf{A different loss picks a different summary.} The same
  population question for the \emph{advantage divisor} $\hat\sigma$ has a
  different answer, because the tilt $\pi e^{\beta Q/\hat\sigma}$ is an
  exponential family and its KL cost is the Bregman divergence of the $Q$
  cumulant generating function: exactly $\tfrac{\beta^2}{2}(\sigma_Q/\hat\sigma-1)^2$
  for Gaussian $Q$ (Proposition~\ref{prop:advkl}). Only the \emph{within}-state
  spread of $Q$ enters, the across-state spread being pure gauge; the cost is
  asymmetric, so one must err high; and the population minimizer is the
  \emph{contraharmonic} mean $\E[\sigma_Q^2]/\E[\sigma_Q]$ --- two accumulators in
  linear space, debiased by $c_4$ (Corollary~\ref{cor:advaccum}), with floor
  $\tfrac{\beta^2}{2}\mathrm{CV}^2/(1+\mathrm{CV}^2)$. Same method, opposite space:
  Remark~\ref{rem:twoscalars}.
  \item \textbf{Nothing depends on $\mu$}, because a translation is known and
  invertible and because mean and covariance propagate in decoupled recursions.
  DDIM transports the mean exactly; for Gaussian data its only error is the loss
  of spread.
\end{enumerate}

\begin{thebibliography}{9}
\bibitem{guo2005mutual}
D.~Guo, S.~Shamai, and S.~Verd\'u.
\newblock Mutual information and minimum mean-square error in Gaussian channels.
\newblock \emph{IEEE Trans.\ Inform.\ Theory}, 51(4):1261--1282, 2005.

\bibitem{nichol2021improved}
A.~Nichol and P.~Dhariwal.
\newblock Improved denoising diffusion probabilistic models.
\newblock \emph{ICML}, 2021. \texttt{arXiv:2102.09672}.

\bibitem{chen2023importance}
T.~Chen.
\newblock On the importance of noise scheduling for diffusion models.
\newblock 2023. \texttt{arXiv:2301.10972}.

\bibitem{karras2022elucidating}
T.~Karras, M.~Aittala, T.~Aila, and S.~Laine.
\newblock Elucidating the design space of diffusion-based generative models.
\newblock \emph{NeurIPS}, 2022. \texttt{arXiv:2206.00364}.

\bibitem{kingma2021vdm}
D.~P.~Kingma, T.~Salimans, B.~Poole, and J.~Ho.
\newblock Variational diffusion models.
\newblock \emph{NeurIPS}, 2021. \texttt{arXiv:2107.00630}.

\bibitem{song2021ddim}
J.~Song, C.~Meng, and S.~Ermon.
\newblock Denoising diffusion implicit models.
\newblock \emph{ICLR}, 2021. \texttt{arXiv:2010.02502}.

\bibitem{chi2023diffusion}
C.~Chi, S.~Feng, Y.~Du, Z.~Xu, E.~Cousineau, B.~Burchfiel, and S.~Song.
\newblock Diffusion policy: visuomotor policy learning via action diffusion.
\newblock \emph{RSS}, 2023. \texttt{arXiv:2303.04137}.

\bibitem{ren2024dppo}
A.~Z.~Ren, J.~Lidard, L.~L.~Ankile, A.~Simeonov, P.~Agrawal, A.~Majumdar,
B.~Burchfiel, H.~Dai, and M.~Simchowitz.
\newblock Diffusion policy policy optimization.
\newblock 2024. \texttt{arXiv:2409.00588}.
\end{thebibliography}

\end{document}
